\chapter{异构多机器人协同}
\label{ch:mr-hetero}

% ════════════════════════════════════════════════════════════════
%  全吸收章：重渲染 Tutorial 笔记
%   - 04_移动机器人规控/50_多机器人协作/70_异构多机协同.md
%     （异构性三轴定义 / 能力互补矩阵 / 异构 CBBA / 空地通信架构 /
%      异构编队阻抗引导 / 跨模态融合 / 异构 wrench cone / HAPPO 异构动作空间）
%  记号归一：R in SO(3)、T in SE(3)（preamble 宏 \SO \SE \so \se）；
%            多机/分配/感知/学习领域自然记号保留，字母隔离就地声明。
%  源由 AI 撰写，论文归属/年份/会议/数值/前沿编号存疑处就地 \pz/\rebuilt，并入 claims 台账。
%  教学层遵循 v5.0 理论教学规范（动机→反面→历史→理论→陷阱→练习 + 误解汇总/小结/速查/故障排查）。
% ════════════════════════════════════════════════════════════════

\rebuilt{本章重渲染自 Tutorial 笔记《异构多机协同——能力互补、空地联合与跨模态融合》，其源稿由 AI 撰写，论文归属（Gerkey--Matari\'c、Choi--Brunet--How、Ravichandar 等）、年份、会议卷期、实验数值（如 RoCo 成功率 86.7\%）与 2024/2025 前沿工作的编号可能有误。本章忠实重述、不擅自“修正”，逐处 \texttt{\textbackslash cite} 标源；存疑断言已登记 claims 台账，待联网核对后二次校验。}

\begin{practice}[前置自测]
开始之前，先试答下列问题。本章每一节都建立在多机协同基础之上，只把“同构”假设替换成“异构”，不重新教共识、CBBA、MPC、grasp matrix 这些底层工具。若答不出 $\ge 2$ 题，建议先回前置章节补齐。
\begin{enumerate}
  \item CBBA 的代价 $c_{ij}$ 表示什么？若某任务只有会飞的机器人才能做，把代价设成很大的正数与设成 $-\infty$（或从可行集剔除）有何本质区别？
  \item 共识 $\dot{\mathbf{x}}=-\mathbf{L}\mathbf{x}$ 收敛到什么？它要求所有 agent 维护同维同义的状态 $x_i$。若无人机维护 $4$ 维位姿、四足维护 $12$ 维关节状态，这个协议还能直接跑吗？问题出在哪？
  \item 分布式 MPC 编队靠什么把各单体的局部优化耦合起来？若一个 agent 是二阶积分点质量、另一个是欠驱动四旋翼（不能瞬间侧移），邻居一致性约束还能写成简单的相对位置相等吗？
  \item grasp matrix $\mathbf{G}$ 把各接触点的力映射到物体合力旋量。若一只手是吸盘（只能拉）、一只手是夹爪（能推能拉）、一架无人机用绳索悬挂（只能沿绳拉），它们的力可行集一样吗？这对内力分配意味着什么？
  \item 四旋翼的微分平坦让规划在 $4$ 维平坦输出空间里做；差速驱动车（unicycle）有“不能横移”的非完整约束。两者的可达运动集在结构上有何根本不同？为什么不能用同一条参考轨迹同时喂给二者？
\end{enumerate}
\noindent\emph{答案线索}：1$\to$\cref{sec:hmr-cbba}（硬剔除 vs 软惩罚）；2$\to$\cref{sec:hmr-formation}（公共输出空间）；3$\to$\cref{sec:hmr-formation}（输出空间一致性）；4$\to$\cref{sec:hmr-manip}（不对称 wrench cone）；5$\to$\cref{sec:hmr-formation,sec:hmr-axes}（可达集结构差异）。
\end{practice}

\section*{本章目标}
学完本章，你应当能够：
\begin{enumerate}
  \item \textbf{形式化异构性}：用“能力 / 动力学 / 感知”三轴拆解一支机器人队伍的异构性，写出能力矩阵 $Q$ 与能力互补矩阵 $C^{\mathrm{comp}}$，并解释可累加特质（cumulative trait）与不可累加特质（non-cumulative trait）的区别为何决定任务能否由联盟拼出；
  \item \textbf{从同构 CBBA 推导异构 CBBA}：把 CBBA 扩展到能力硬约束、联盟任务与能量约束三种异构因素，论证“硬剔除”优于“软惩罚”，并解释引入联盟任务为何可能破坏边际收益递减（DMG）条件、$50\%$ 最优界因此失效；
  \item \textbf{设计空地联合系统}：讲清空地异构系统（AGHS）的三类通信架构（中继 / mesh / leader--follower）各自的单点故障特性，并能为搜救、巡检、物流选对架构；
  \item \textbf{实现跨模态信息融合}：说明视角鸿沟（viewpoint gap）与模态鸿沟（modality gap）为何让简单位姿广播失效，掌握 $2.5\mathrm{D}$ 高程图配准、深度描述子检索、半分布式相对定位三条路线；
  \item \textbf{构建异构编队}：理解异构编队的一致性为何必须建在公共输出空间而非状态空间，掌握 leader--follower $+$ 阻抗引导（impedance-based guidance）如何让地面车忽略“无人机能飞过但地面车要绕开”的矮障碍；
  \item \textbf{推广内力分配到异构接触}：把对称摩擦锥下的内力均分推广到不对称 wrench cone 上的带约束优化，并用 force closure 判断异构抓取可行性；
  \item \textbf{理解异构 MARL}：解释 HAPPO 如何用多智能体优势分解 $+$ 顺序更新处理维度/语义不一致的动作空间，说明它为何必须用无参数共享（NoPS）才能保证单调改进、以及由此带来的样本效率代价；
  \item \textbf{辨析异构与同构的边界}：用一张五层对比总表说清“异构系统不是同构系统的并集”，并能为真实“四足 $+$ 无人机”场景判断哪些工具可直接复用、哪些必须改写。
\end{enumerate}

\subsection*{本章知识导航}
\textbf{一句话定位}：前面的多机章节默认所有机器人是“同一种东西”——同样的动力学、传感器、动作空间。本章把这个假设拆掉：当队伍里同时有“看得远的无人机、搬得重的四足、抓得准的机械臂”时，协同问题的每一层都要重写。一条认知主线贯穿全章——\textbf{先定义异构，再看异构如何改写每一层协同，最后归纳异构与同构的本质边界}。

\begin{center}
\small
\begin{tikzpicture}[
  node distance=4mm and 5mm, >=Stealth,
  blk/.style={draw, rounded corners, align=center, font=\small, inner sep=3pt, fill=main!6, draw=main!60!black},
  grp/.style={draw, rounded corners, align=center, font=\small, inner sep=3pt, fill=second!6, draw=second!60!black},
  fin/.style={draw, rounded corners, align=center, font=\small, inner sep=3pt, fill=third!8, draw=third!60!black},
  ln/.style={->, thick, gray!60!black}]
\node[blk] (root) {什么是异构？\\三轴定义 $+$ 能力互补矩阵};
\node[grp, below=9mm of root] (alloc) {异构分配\\能力约束 CBBA $+$ 联盟};
\node[grp, left=of alloc] (comm) {异构通信\\中继 / mesh / LF};
\node[grp, right=of alloc] (form) {异构编队\\阻抗引导};
\node[grp, below=8mm of comm] (perc) {跨模态感知\\视角 / 模态鸿沟};
\node[grp, below=8mm of alloc] (manip) {异构操作\\不对称 wrench cone};
\node[grp, below=8mm of form] (learn) {异构学习\\HAPPO 动作空间};
\node[fin, below=8mm of manip] (bound) {异构 vs 同构边界\\五层总表 $+$ 寻找公共空间};
\draw[ln] (root)--(alloc); \draw[ln] (root.west) to[out=200,in=90] (comm.north);
\draw[ln] (root.east) to[out=-20,in=90] (form.north);
\draw[ln] (comm)--(perc); \draw[ln] (alloc)--(manip); \draw[ln] (form)--(learn);
\draw[ln] (perc.south) |- (bound.west); \draw[ln] (manip)--(bound);
\draw[ln] (learn.south) |- (bound.east);
\end{tikzpicture}
\end{center}

\noindent\textbf{推荐路径}：\cref{sec:hmr-axes}（三轴定义 $+$ 能力互补矩阵）是全章地基与共同语言，\emph{必读}；\cref{sec:hmr-cbba}（异构 CBBA）把分配异构化，是分配线的核心；\cref{sec:hmr-comm}（通信架构）偏系统、相对轻松；\cref{sec:hmr-formation}（异构编队）与 \cref{sec:hmr-perception}（跨模态融合）、\cref{sec:hmr-marl}（异构 MARL）是三个硬核小节——其中 \cref{sec:hmr-perception}、\cref{sec:hmr-marl} 是本章真正区别于“同构多机”的地方，值得花最多时间；\cref{sec:hmr-manip}（异构操作）把内力分配异构化；\cref{sec:hmr-bound}（本质边界）是认知收口。两条阅读线：\textbf{系统线}（搭异构系统）\cref{sec:hmr-axes}$\to$\cref{sec:hmr-cbba}$\to$\cref{sec:hmr-comm}$\to$\cref{sec:hmr-formation}；\textbf{算法线}（吃透新数学）\cref{sec:hmr-axes}$\to$\cref{sec:hmr-cbba}$\to$\cref{sec:hmr-perception}$\to$\cref{sec:hmr-marl}$\to$\cref{sec:hmr-bound}。

\subsection*{前置知识桥接}
本章紧接多机基础（图论、共识）与协同运动（分布式 MPC、协同搬运），这里把最关键的前置点重新激活——你不必翻回去就能跟上：
\begin{itemize}
  \item \textbf{共识作为分布式裁决的原子操作}（\cref{ch:mr-consensus}）。共识 $\dot{\mathbf{x}}=-\mathbf{L}\mathbf{x}$ 让各 agent 只靠邻居通信收敛到一致（双随机权重下收敛到平均）。本章会两次用到它，但都要先做一步“投影”：\cref{sec:hmr-cbba} 的异构 CBBA 冲突消解仍用标量最大值共识（异构不影响——出价是标量），而 \cref{sec:hmr-formation} 的几何一致性必须先把异构状态投影到公共位置空间再做共识。\textbf{异构性不改变共识的数学，但改变了“对什么做共识”}。
  \item \textbf{CBBA $=$ 分布式拍卖 $+$ 共识}。CBBA 分两阶段：本地 bundle 构建（贪婪地往任务包里加边际收益最高的任务）$+$ 基于共识的冲突消解（带时间戳的最大值共识裁决谁中标），其 $50\%$ 最优界依赖边际收益递减（DMG）条件。本章 \cref{sec:hmr-cbba} 在这个骨架上加三样东西——能力硬约束、联盟任务、能量约束；\textbf{两阶段结构原封不动，变的是喂给它的问题定义}。
  \item \textbf{分布式 MPC 编队}（\cref{ch:mr-consensus} 的编队范式）。一致性约束写在相对位置上、各机器人共享同一动力学模板。本章 \cref{sec:hmr-formation} 指出：动力学异构时一致性不能再写在状态空间（不同 agent 状态不同维），必须上移到输出空间（编队几何）。
  \item \textbf{grasp matrix 与内力分配}。grasp matrix $\mathbf{G}$ 把接触力映射到物体合力旋量，内力落在 $\mathbf{G}$ 的零空间里；同构抓取默认各 agent 力可行集对称、内力均分。本章 \cref{sec:hmr-manip} 指出：异构接触（吸盘 / 夹爪 / 绳索）的力可行集形状各异，内力分配要在不对称的 wrench cone 上做投影。
\end{itemize}

\subsection*{如果跳过本章会怎样}
跳过本章，你会在三个具体的地方栽跟头。
\begin{itemize}
  \item \textbf{“我把 CBBA 代价设成大正数，结果无人机被派去搬重物了”}：为表达“无人机搬不动”把对应 $c_{ij}$ 设成 $10^6$，大多数时候没问题，但某次实例里所有地面车都满载，优化器为“完成所有任务”硬把搬运派给无人机——因为 $10^6$ 虽大却有限。\cref{sec:hmr-cbba} 会告诉你：能力约束必须是硬剔除而非软惩罚。
  \item \textbf{“无人机和地面车各自定位很准，拼一起却差两米”}：各自跑 SLAM 再广播位姿想拼成共享地图，结果两图错位严重——鸟瞰特征（屋顶、道路）和第一视角特征（墙面、家具）几乎无共视，回环失效。\cref{sec:hmr-perception} 会告诉你：跨模态、跨视角存在视角鸿沟和模态鸿沟，简单位姿广播解决不了坐标系对齐。
  \item \textbf{“我想用一个 MAPPO 同时训无人机和四足，动作维度对不上直接报错”}：无人机 $4$ 维控制、四足 $12$ 维关节，参数共享网络搭不起来；padding 到同维后训练完全不收敛。\cref{sec:hmr-marl} 会告诉你：异构动作空间是 MARL 的真问题，HAPPO 用顺序更新 $+$ 独立策略处理，但代价是放弃参数共享、样本效率下降。
\end{itemize}

\begin{table}[H]\centering\small
\caption{本章阅读方式建议}
\label{tab:hmr-reading}
\begin{tabularx}{\linewidth}{LlL}
\toprule
阅读方式 & 时间 & 适合谁 \\
\midrule
精读（逐节读动机→推导→例，亲手实现异构 CBBA / 编队 / 配准） & 14--18 小时 & 第一次系统学异构多机协同 \\
速读（读每节动机与理论、能力矩阵、五层总表、端到端案例） & 4--6 小时 & 已有同构多机背景，建全局图景 \\
速查（能力矩阵 $+$ 异构 CBBA 三改动 $+$ 通信对比 $+$ 融合三路线 $+$ HAPPO $+$ 五层总表） & 50--70 分钟 & 已学过，回查特定内容 \\
\bottomrule
\end{tabularx}
\end{table}

\begin{note}
\textbf{本章记号与约定.}\;
旋转矩阵记 $\mathbf{R}\in\SO(3)$、位姿 $\mathbf{T}\in\SE(3)$（preamble 宏 $\SO,\SE,\so,\se$，与全书一致，\nameref{ch:notation}）；多机/分配/感知/学习领域的自然记号保留。
机器人集合 $\mathcal{R}=\{r_1,\dots,r_N\}$，可分多个物种（species）$\mathcal{S}=\{s_1,\dots,s_K\}$，物种内同构、物种间异构。
机器人 $i$ 的能力/特质向量 $\mathbf{q}_i\in\mathbb{R}^U$（$U$ 维），能力矩阵 $Q\in\mathbb{R}^{N\times U}$（行 $=$ 机器人、列 $=$ 能力维度）；任务 $j$ 的能力需求向量 $\mathbf{y}_j\in\mathbb{R}^U$；$\succeq$ 表逐分量大于等于。
能力互补矩阵 $C^{\mathrm{comp}}$；合格候选集 $\mathcal{A}_j$；机器人 $i$ 在接触点的可施加 wrench 锥 $\mathcal{W}_i$。
\textbf{字母隔离声明}：本章 $c_{ij}$ 指分配代价、$\alpha/\beta$ 在 \cref{sec:hmr-cbba} 指能量代价权重、在 \cref{sec:hmr-formation} 阻抗系数处就地另释；$A^{i_{1:m}}$ 指 HAPPO 多智能体联合优势、$\pi^i$ 指 agent $i$ 的独立策略、$M$ 指顺序更新的累积重要性比率，均仅本章局部生效。
\end{note}

% ════════════════════════════════════════════════════════════════
\section{异构性的三轴定义与能力互补矩阵}
\label{sec:hmr-axes}

\paragraph{动机：一个“显然”却没人说清的词。}
打开任何一篇多机器人论文，“异构（heterogeneous）”出现频率极高，但几乎没有两篇用同一个定义：有人指型号不同、有人指动力学不同、有人指传感器配置不同、还有人指跑不同算法。这种含糊在写综述时无害，但当你要\emph{实现}一个异构协同算法时，它立刻变成具体障碍——代码里描述“机器人 $i$ 能做什么”的数据结构该长什么样？是枚举类型（UAV/UGV/ARM）、能力布尔向量，还是连续数值向量？这三种选择对应三种完全不同的算法设计。更糟的是，含糊定义会让你\textbf{误判难度}：一个常见错误是“反正都是机器人，把同构算法的参数改一改就行”——这对某些异构维度成立（最大速度不同，确实只是参数），对另一些完全不成立（动作空间维度不同，参数共享网络根本搭不起来）。\textbf{不区分“哪种异构只是参数差异、哪种是结构差异”，是异构系统设计中最致命的认知错误。}

\paragraph{反面：用最朴素的“型号标签”会怎样。}
最朴素的表示是给每个机器人贴型号标签（无人机/四足/机械臂），代码里就是枚举。它有三个致命缺陷：（1）\textbf{标签太粗，丢失能力连续性}——载重 $0.2\,\mathrm{kg}$ 的微型四旋翼和载重 $5\,\mathrm{kg}$ 的工业六旋翼共享“无人机”标签，分配器无法区分“这个搬运需要 $3\,\mathrm{kg}$ 载重，微型机做不了”；（2）\textbf{标签是封闭集，无法表达能力的组合与缺失}——同一台四足装了机械臂就能操作、没装就不能，能力实际是“型号 $+$ 配置 $\Rightarrow$ 能力”而非“型号 $\Rightarrow$ 能力”；（3）\textbf{标签无法支持能力累加的推理}——异构最有价值的场景是联盟（coalition），单机做不到的任务由多机能力拼合，但“搬 $8\,\mathrm{kg}$”可由两台 $5\,\mathrm{kg}$ 机器人完成（载重可累加），“飞到 $120\,\mathrm{m}$”却不能由两台 $60\,\mathrm{m}$ 无人机拼出（高度不可累加），型号标签完全无法支持这种推理。

\begin{insight}[以能力为中心，而非以型号为中心]
型号标签的根本错误在于它把“机器人是什么”和“机器人能做什么”混为一谈。异构协同关心的从来不是“它是什么型号”，而是“它具备哪些能力、这些能力能否与他人累加”。\textbf{正确的异构表示应以能力（capability/trait）为中心}——这正是 STRATA 框架\cite{ravichandar2019strata} 的核心洞见，也是本节能力矩阵表示的出发点。
\end{insight}

\paragraph{历史：从型号分类到特质向量。}
异构多机的表示经历了三代演进。\emph{第一代——型号分类}（2000s 早期）：按物理型号组织，分配用规则匹配（“侦察任务 $\to$ 派无人机”），简单但有上述三缺陷。\emph{第二代——能力分类法}（Gerkey \& Matari\'c, IJRR 2004\cite{gerkey2004taxonomy}）：多机任务分配（MRTA）分类法第一次系统地把“任务对能力的需求”形式化，三维分类——单任务/多任务机器人（ST/MT）、单机器人/多机器人任务（SR/MR）、瞬时/时间扩展分配（IA/TA）——让异构能力有了讨论框架，其中 SR/MR 直接刻画“任务能否单机完成”，是异构联盟的雏形；但这一代仍把能力当作离散的“有/无”。\emph{第三代——连续特质向量}（STRATA, Ravichandar 等\cite{ravichandar2019strata}）\rebuilt{源标注会议为 ICRA 2019，经核有误：STRATA 正式发表于 JAAMAS 2020（doi:10.1007/s10458-020-09461-y），亦见 AAMAS 2021；arXiv:1903.05149（2019）仅为预印本，ICRA 2019 只是 workshop 展示。论文实质描述无误，仅会议归属/bib 年份应改为 JAAMAS 2020}：做了关键概念跃迁，把每个机器人的能力建模为连续特质向量（trait vector），每一维是可量化的能力强度（载重 $\mathrm{kg}$、续航 $\min$、最高速度、传感器分辨率……），并区分两类特质：
\begin{itemize}
  \item \textbf{可累加特质（cumulative trait）}：多机的该能力可\emph{相加}满足任务需求，如总载重（两台 $5\,\mathrm{kg}$ 合搬 $10\,\mathrm{kg}$）；
  \item \textbf{不可累加特质（non-cumulative trait）}：该能力\emph{不能}通过多机相加满足，如最高飞行高度（两台 $60\,\mathrm{m}$ 加起来仍只能到 $60\,\mathrm{m}$）、最高分辨率、最快单点速度。
\end{itemize}
STRATA 还借用生物多样性概念，把机器人按特质聚成\textbf{物种（species）}——物种内同构、物种间异构，给出干净的两层结构，本章后面所有算法都建立在此之上。

\begin{note}
本节的“三轴定义”是在 STRATA 特质向量基础上的\emph{教学性扩展}\rebuilt{（三轴划分为本书重述，非 STRATA 原文术语，待核）}。STRATA 主要关注任务分配层的能力特质；本章把异构性扩展到三个轴（能力 / 动力学 / 感知），因为异构不只影响分配，还影响编队（动力学轴，\cref{sec:hmr-formation}）和感知融合（感知轴，\cref{sec:hmr-perception}）。
\end{note}

\subsection{异构性的三轴定义}
\label{sec:hmr-three-axes}
把一支机器人队伍的异构性沿三个\emph{正交}的轴刻画。之所以是三轴而非一个笼统的“异构度”，是因为\textbf{不同轴的异构给不同层的协同带来不同麻烦}——混在一起谈，就会犯“以为改改参数就行”的错误。

\paragraph{轴一：能力异构——影响“谁干什么”。}
指机器人\emph{能完成的任务种类与强度}不同。这是任务分配（\cref{sec:hmr-cbba}）关心的唯一轴。用能力向量刻画每个机器人：
\begin{equation}
  \mathbf{q}_i = (q_{i,1},q_{i,2},\dots,q_{i,U})^\top \in \mathbb{R}^U,
  \label{eq:hmr-qi}
\end{equation}
$q_{i,u}$ 是机器人 $i$ 在第 $u$ 种能力上的强度。按行堆叠得\textbf{能力矩阵}
\begin{equation}
  Q = \begin{pmatrix}\mathbf{q}_1^\top\\\mathbf{q}_2^\top\\\vdots\\\mathbf{q}_N^\top\end{pmatrix}\in\mathbb{R}^{N\times U},\qquad Q_{iu}=q_{i,u},
  \label{eq:hmr-Q}
\end{equation}
行是机器人、列是能力维度。这个矩阵是本章所有分配算法的输入。以一支“$2$ 无人机 $+$ $2$ 四足 $+$ $1$ 机械臂”的队伍为例，取 $U=4$ 个能力维度（载重 $\mathrm{kg}$ / 鸟瞰感知范围 $\mathrm{m}$ / 操作精度 $\mathrm{mm}^{-1}$ / 越障通过性 $0$--$1$）：
\begin{equation}
  Q = \begin{pmatrix}
  0.5 & 150 & 0 & 1.0 \\
  0.5 & 150 & 0 & 1.0 \\
  30 & 20 & 2 & 0.8 \\
  30 & 20 & 2 & 0.8 \\
  50 & 5 & 100 & 0
  \end{pmatrix}
  \begin{array}{l}
  \leftarrow \mathrm{UAV}_1 \\ \leftarrow \mathrm{UAV}_2 \\ \leftarrow \text{四足}_1 \\ \leftarrow \text{四足}_2 \\ \leftarrow \text{臂}
  \end{array}
  \label{eq:hmr-Q-example}
\end{equation}
读这个矩阵就能看出能力互补的结构：无人机的感知范围列（$150$）远高于其他，机械臂的操作精度列（$100$）一骑绝尘，四足在载重和通过性上居中且均衡。\textbf{异构性在这里变成了可计算、可比较、可优化的对象}。每个任务也用能力需求向量 $\mathbf{y}_j\in\mathbb{R}^U$ 描述，语义是“完成任务 $j$ 至少需要在每个能力维度上达到 $\mathbf{y}_j$ 的强度”，任务-能力匹配退化为向量比较（\cref{sec:hmr-cbba} 详述）。

\begin{insight}[互补 $=$ 能力矩阵列覆盖性的定量性质]
把能力写成向量、把队伍写成矩阵，最深刻的好处不是记法简洁，而是\textbf{它让“互补”这个直觉概念变得可度量}。两台机器人互补，数学上意味着它们的能力向量在不同维度上有峰值（“你强我弱、我强你弱”）；一支队伍能力全面，意味着能力矩阵的每一列都至少有一个大值。互补不再是定性的赞美，而是能力矩阵列覆盖性的定量性质。
\end{insight}

\paragraph{轴二：动力学异构——影响“怎么动”。}
指机器人的\emph{运动模型在结构上}不同——不是参数不同（那只是同一模型的不同实例），而是\emph{模型方程的形式、状态维度、约束类型}不同。这是编队（\cref{sec:hmr-formation}）关心的轴。举三个结构上根本不同的运动模型（\cref{tab:hmr-dynamics}）：
\begin{table}[H]\centering\small
\caption{三类结构不同的运动模型}
\label{tab:hmr-dynamics}
\begin{tabular}{llll}
\toprule
机体 & 运动模型 & 状态维度 & 关键约束 \\
\midrule
二阶积分点质量 & $\ddot{p}=u$ & $6$（位置 $+$ 速度） & 无（全向） \\
差速驱动车（unicycle） & $\dot x=v\cos\theta,\ \dot y=v\sin\theta,\ \dot\theta=\omega$ & $3$ & 非完整：不能横移 \\
四旋翼（欠驱动） & $\SE(3)$ 刚体 $+$ 单向推力 & $12$ & 欠驱动：水平加速靠倾斜 \\
\bottomrule
\end{tabular}
\end{table}
\noindent 三者的差异不是“最大速度不同”这种参数差异，而是\textbf{可达运动集（reachable motion set）的结构差异}：点质量的可达瞬时速度方向各向同性；差速车被锁在车头朝向（一维，亏秩）；四旋翼位置上各向同性（微分平坦），但姿态-位置强耦合且推力单向。直接后果是：\textbf{一条对某类机体可行的轨迹，喂给另一类可能直接违反约束}——含横向平移段的轨迹对四旋翼可行（侧飞靠倾斜），对差速车不可行（非完整约束禁止横移）。这就是为什么 \cref{sec:hmr-formation} 说异构编队的一致性必须建在\emph{输出空间}而非状态空间。

\begin{pitfall}[对比性思维：动力学异构不是“同模型不同参数”，而是“不同结构的模型”]
参数异构可用同一套控制器加\emph{增益调度}（gain scheduling）处理——这其实仍是同构控制的范畴；结构异构必须为每类机体设计各自的控制器，再在更高层协调。把结构异构误当参数异构，是异构编队失败的头号原因——你会发现一套精心调好的控制器在另一类机体上怎么都不稳定，因为它假设的运动模型根本不对。
\end{pitfall}

\paragraph{轴三：感知异构——影响“看到什么”。}
指机器人的\emph{传感模态、视角、坐标系}不同。这是跨模态融合（\cref{sec:hmr-perception}）关心的轴，可细分为三个子维度，每个都让简单的“位姿广播 $+$ 地图拼接”失效：\textbf{模态异构}（无人机下视相机、地面车激光、搜救机器人热成像——图像 vs 点云 vs 温度场，无法直接比对）；\textbf{视角异构}（鸟瞰 vs 平视——同一物体在两个视角下可能完全无共视特征，屋顶和墙面是同一栋楼的不同面）；\textbf{坐标系异构}（每个机器人有自己的本体系 $B_i$ 和里程计原点，各自 SLAM 的地图在各自坐标系里，无共同世界系基准）。三者叠加构成 \cref{sec:hmr-perception} 要攻克的视角鸿沟 $+$ 模态鸿沟。

\begin{insight}[异构既是负担又是红利]
三轴异构可从两个互补角度理解，二者指向同一工程结论。\emph{负担视角}：三个轴各自破坏了同构系统赖以工作的某个假设（能力轴破坏“谁都能干”、动力学轴破坏“状态可比”、感知轴破坏“观测可拼”），所以异构系统每一层都要额外做“消解差异”的工作。\emph{资源视角}：正因能力轴有差异才有互补（无人机看得远、四足搬得重）；正因动力学轴有差异才能覆盖更广运动模式（飞越 $+$ 越障）；正因感知轴有差异才能多模态互证（鸟瞰看全局 $+$ 平视看细节）。\textbf{工程的艺术就是用最小的“消解差异”代价换取最大的“能力互补”收益}——本章每一节本质上都在做这个权衡。
\end{insight}

\paragraph{三轴的正交性与耦合。}
这三个轴\emph{概念上正交}——一个机器人可以只在某一轴上异构（两台同型号无人机一台装相机一台装激光，则能力/动力学轴同构、仅感知轴异构）。但真实系统里三轴往往\emph{同时}异构（无人机 vs 四足三轴全异构），这正是异构协同难的根源：分配把任务派给无人机（能力轴决策），但无人机的轨迹必须由它自己的动力学生成（动力学轴），而它感知到的目标位置又在它自己的坐标系里（感知轴），三者必须打通才能让任务落地（\cref{sec:hmr-bound} 的端到端案例演示这三轴如何在一个系统里协同打通）。

\subsection{能力互补矩阵：把“互补”变成可计算的对象}
\label{sec:hmr-comp}
有了能力矩阵 $Q$，可进一步定义\textbf{能力互补矩阵} $C^{\mathrm{comp}}$，度量队伍中任意两机器人的互补程度。一个 pair $(i,k)$ 互补的直觉是：在某些维度上 $i$ 强 $k$ 弱、另一些维度上 $k$ 强 $i$ 弱。先把每个能力向量按列归一化（消除量纲差异，载重 $\mathrm{kg}$ 与高度 $\mathrm{m}$ 不能直接比）得 $\hat{\mathbf{q}}_i$，互补度定义为
\begin{equation}
  C^{\mathrm{comp}}_{ik} = \frac{1}{U}\sum_{u=1}^{U}\big|\hat q_{i,u}-\hat q_{k,u}\big|\cdot\mathbb{1}\big[\text{维度 }u\text{ 上两者有显著差异}\big].
  \label{eq:hmr-comp}
\end{equation}
$C^{\mathrm{comp}}_{ik}$ 越大表示 $i,k$ 越互补（差异越大、越能填补彼此短板），越小表示越同质（能力雷同、放在一起冗余）。注意它与“相似度”恰好相反。

\begin{insight}[同构是“互补度处处为零”的极限情形]
如果一支队伍的 $C^{\mathrm{comp}}$ 全是接近 $0$ 的值，意味着所有机器人能力雷同——这恰恰\emph{退化回了同构系统}。反过来，互补矩阵的值越大、结构越丰富，队伍的异构红利越大。这给了组建异构队伍的定量指南：\textbf{好的异构队伍不是机型越多越好，而是能力覆盖越全、互补结构越互不重叠越好}。
\end{insight}

\paragraph{联盟能力聚合：区分可累加性。}
异构联盟分配的核心运算是“多个异构机器人合起来能做什么”，归约为一个\emph{区分可累加性}的向量聚合。给定联盟成员集 $S$，其聚合能力第 $u$ 维为
\begin{equation}
  \mathrm{agg}_u(S) =
  \begin{cases}
  \sum_{i\in S} q_{i,u}, & u\text{ 为可累加维（如总载重）},\\[4pt]
  \max_{i\in S} q_{i,u}, & u\text{ 为不可累加维（如最高视野）}. % [待修] 此 max 非 STRATA 原文(原文为阈值计数)，见下方 \rebuilt 注；保留为本书工程约定
  \end{cases}
  \label{eq:hmr-agg}
\end{equation}
STRATA 框架的确形式化了 cumulative / non-cumulative 这个区分。\rebuilt{但 \cref{eq:hmr-agg} 的“不可累加维取 $\max$”并非 STRATA 原文：经核对原文（arXiv:1903.05149 式(13)(17)），STRATA 对所有特质都用同一个线性求和聚合 $Y=XQ$，不可累加维的处理是“先按最小需求二值化、再计数（sum of indicators）”，即 threshold-then-count，而非 $\max$。此处 $\max$ 为本书的工程化简写、可作独立设计取用，但勿当作 STRATA 公式。正确方向：若严格复刻 STRATA，应改为阈值计数；可累加维取 sum 的一支与原文方向一致。}本式作为本书联盟聚合的工程约定，\cref{sec:hmr-cbba} 的联盟任务会反复用到它。下面给出能力矩阵与互补矩阵的构造，作为本章累积项目的第一块拼图（注意能力维度的语义、可累加性标注、行列方向都要固定，避免后续串味）。

\begin{codebox}[Python]{能力矩阵、互补矩阵与联盟能力聚合}
import numpy as np

CAP_NAMES  = ["payload_kg", "aerial_range_m", "precision_inv_mm", "traversability"]
CUMULATIVE = np.array([True, False, False, False])   # payload 可累加;其余不可累加
Q = np.array([                                       # 行=机器人，列=能力维度
    [ 0.5, 150,   0, 1.0],   # UAV_1
    [ 0.5, 150,   0, 1.0],   # UAV_2
    [30.0,  20,   2, 0.8],   # quadruped_1
    [30.0,  20,   2, 0.8],   # quadruped_2
    [50.0,   5, 100, 0.0],   # arm
], dtype=float)

def normalize_columns(Q):                            # 按列归一化到 [0,1]，消除量纲差异
    col_min, col_max = Q.min(0), Q.max(0)
    span = col_max - col_min; span[span == 0] = 1.0  # 常数列防除零
    Qn = (Q - col_min) / span
    Qn[:, (col_max - col_min) == 0] = 0.5            # 常数列无信息取中性值
    return Qn

def complementarity_matrix(Q, sig=0.3):              # C[i,k]: i,k 各维显著差异均值，越大越互补
    Qn = normalize_columns(Q); N, U = Qn.shape; C = np.zeros((N, N))
    for i in range(N):
        for k in range(N):
            if i == k: continue
            diff = np.abs(Qn[i] - Qn[k])
            C[i, k] = np.sum(diff * (diff > sig)) / U
    return C

def coalition_capability(Q, members, cumulative=CUMULATIVE):   # 可累加维 sum、不可累加维 max
    sub = Q[members]
    return np.where(cumulative, sub.sum(0), sub.max(0))
\end{codebox}

\noindent 期望观察：UAV 与 Arm 互补度最高（感知 vs 精度两个极端），$\mathrm{UAV}_1$ 与 $\mathrm{UAV}_2$ 互补度为 $0$（完全同质、冗余）；联盟 $\{$Quad$_1$,Quad$_2\}$ 的聚合载重为 $30+30=60$（可累加），鸟瞰范围为 $\max(20,20)=20$（不可累加——视野不会因人多而变远）。常见错误有三：\emph{不归一化就算差异}（互补度被量级最大的列主导）；\emph{把不可累加维当可累加}（“两台 $150\,\mathrm{m}$ 视野无人机合起来 $300\,\mathrm{m}$”物理荒谬）；\emph{行列方向搞反}（$N\ne U$ 时才报错、$N=U$ 时静默出错，须 \texttt{assert Q.\allowbreak shape[0]==len(ROBOT\_NAMES)}）。

\begin{pitfall}[编程陷阱：不可累加能力被错误累加]
判断联盟能否完成任务时对所有能力维度一律求和，会让算法认为“两台最高飞 $60\,\mathrm{m}$ 的无人机组队能飞到 $120\,\mathrm{m}$”，把“飞到 $100\,\mathrm{m}$”的任务派给联盟，执行时永远飞不到、任务永久失败。\textbf{根本原因}：把不可累加特质（最高高度/分辨率/单点速度）当成了可累加特质——可累加性是能力的内在属性。\textbf{正确做法}：为每维标注 cumulative 布尔位，聚合时可累加维 \texttt{sum}、不可累加维 \texttt{max}（\cref{eq:hmr-agg}）。\textbf{自检}：对每维问“两个机器人一起做，这个指标会翻倍吗？”——载重会，最高速度/高度/精度都不会。
\end{pitfall}

\begin{pitfall}[概念误区：“机型越多，异构系统能力越强”]
往队伍里多加机型，系统未必更强。\textbf{现象}：加了第三、四种机型，复杂度（通信、协调、维护）暴涨，但能完成的任务种类没相应增加——因为新机型能力与已有机型重叠。\textbf{根本原因}：异构红利来自能力的“覆盖性”和“互补性”，不是机型数量；若新机型能力向量与已有高度相似（互补矩阵该行该列接近 $0$），它只增冗余、不扩边界。\textbf{正确理解}：判断“该不该加这种机型”要看它是否填补能力矩阵某一列的空缺，或在某关键能力上提供更高强度。
\end{pitfall}

\begin{pitfall}[思维陷阱：把“参数异构”当成“结构异构”，或反之]
\textbf{错误之一}：看到两机器人最大速度不同就为它们设计完全不同的协同协议——过度设计（本来加增益调度即可）。\textbf{错误之二}：看到两机器人都“会移动”就用同一条含横移段轨迹喂给差速车和四旋翼——欠设计（差速车直接违反非完整约束）。\textbf{根本原因}：没区分参数异构（同模型不同参数，如 $\max v$）与结构异构（不同模型，如全向 vs 非完整 vs 欠驱动）。\textbf{正确做法}：面对异构队伍，先沿三轴判断每个轴是参数还是结构异构——能力轴几乎总是强度差异（统一在能力矩阵里处理），动力学轴看运动模型方程是否同形，感知轴看模态/视角/坐标系是否同类。
\end{pitfall}

\begin{practice}
\begin{enumerate}
  \item （实现 $+$ 分析）用 \texttt{complementarity\_matrix} 计算 \cref{eq:hmr-Q-example} 那支队伍的完整 $5\times5$ 互补矩阵。(a) 哪一对互补度最高、为什么？(b) 若只能保留 $3$ 台以最大化“能力列覆盖”，你保留哪 $3$ 台？讨论“最大化互补度之和”与“最大化能力覆盖”是否一致，构造反例说明它们可能冲突。
  \item （设计题）为“灾后建筑结构评估”设计能力维度（$\ge 5$ 维）与任务需求向量（高空裂缝拍摄、室内梁柱检测、承重测试、有毒气体检测）。对每维标注 cumulative/non-cumulative 并解释依据；说明哪些任务必须由联盟完成、哪些可单机完成。
  \item （跨章综合，串联图论与 CBBA）考虑通信半径异构的队伍（无人机 $500\,\mathrm{m}$、地面车 $100\,\mathrm{m}$）。论证：(a) 异构通信半径如何让通信图 $\mathcal{G}$ 本身变成“异构图”（边的存在取决于哪类节点）；(b) 无人机失效时 $\lambda_2$ 可能如何变化、甚至图不连通；(c) 这对图上运行的 CBBA 收敛性的影响，并联系 \cref{sec:hmr-comm} 的“中继型架构单点故障”。
\end{enumerate}
\end{practice}

% ════════════════════════════════════════════════════════════════
\section{异构任务分配：能力感知的 CBBA 扩展}
\label{sec:hmr-cbba}

\paragraph{动机：当“谁都能做”的假设崩塌。}
CBBA 解决的是干净问题：$N$ 个同质机器人、$M$ 个任务，每个机器人对每个任务有代价 $c_{ij}$，目标是分配使总效用最大；其优雅在完全分布式——本地构建任务包、靠共识消解冲突，拿到不差于最优 $50\%$ 的解\cite{choi2009cbba}。但它有一个隐含的同质假设：\textbf{任何机器人都可以竞标任何任务}。异构队伍里这直接崩塌——搜救场景里“搬开挡路的预制板”需要大载重，让无人机竞标它毫无意义（搬不动）；“航拍废墟全景”需要飞行，让地面车竞标同样荒谬。若天真地把不可行分配的代价设成大数（如 $10^6$），会出现隐蔽而致命的 bug：\textbf{在某些边角情况下优化器为“完成所有任务”会选一个超贵但有限的不可行分配}（所有地面车满载、“搬预制板”无人认领时，优化器认为“派无人机去搬，代价 $10^6$，总比留着任务不完成强”，于是把物理上不可能的任务派给无人机，仿真里表现为无人机收到永远无法完成的指令、卡死）。这引出异构分配第一个、也是最基础的改动：\textbf{能力约束必须是硬的，不可行分配必须从竞标资格里彻底剔除。}

\paragraph{反面：硬约束 vs 软约束。}
\emph{软约束（大代价）}保留所有 $(i,j)$ 对、把不可行的设成 $c_{ij}=M_{\mathrm{big}}$，它在优化目标里仍是可选项，只是“很贵”。\emph{硬约束（剔除）}为每个任务 $j$ 维护\textbf{合格候选集} $\mathcal{A}_j=\{i:\mathbf{q}_i\succeq\mathbf{y}_j\}$，只有 $\mathcal{A}_j$ 里的机器人能竞标 $j$，不在候选集里的 pair 根本不进入优化。两者区别见 \cref{tab:hmr-hard-soft}。

\begin{table}[H]\centering\small
\caption{能力约束：软约束（大代价）vs 硬约束（剔除候选）}
\label{tab:hmr-hard-soft}
\begin{tabular}{lll}
\toprule
层面 & 软约束（大代价） & 硬约束（剔除候选） \\
\midrule
正确性 & 极端情况不可行分配仍可能被选中 & 不可行分配永不出现（结构上保证） \\
数值 & 大代价污染出价，$\varepsilon$-松弛/收敛阈被扭曲 & 出价干净，只在可行集上比较 \\
效率 & 在大量“必不选”项的空间里搜索 & 搜索空间直接缩小到可行子集 \\
\bottomrule
\end{tabular}
\end{table}

\begin{insight}[用约束表达“绝对不可以”，用代价表达“不太好但可以”]
硬约束优于软约束，背后是一个超越异构分配的普适原则——\textbf{用约束的结构（可行集的形状）表达“绝对不可以”，用代价的数值表达“不太好但可以”}。把“绝对不可以”编码成“代价很大”，本质是把一个布尔事实（能/不能）降级成连续量（贵/便宜），而连续量在优化中永远存在“被其他更糟的选择反超”的可能。这个原则在 \cref{sec:hmr-formation} 的避障（碰撞是硬的）、\cref{sec:hmr-manip} 的不对称内力分配（接触约束是硬的）里会反复出现：\textbf{凡是物理上“绝对不可能”的，都应编码进可行集而非代价函数。}
\end{insight}

\paragraph{历史：从同质拍卖到能力感知分配。}
CBBA（Choi, Brunet, How, IEEE T-RO 2009\cite{choi2009cbba}）本身是同质的。把它异构化的需求很早出现——Gerkey--Matari\'c MRTA 分类法（2004）的 SR/MR 维度（单机器人任务 vs 多机器人任务）就预示了“联盟任务”的存在。真正系统地处理异构能力分配的是 STRATA（Ravichandar 等, ICRA 2019\cite{ravichandar2019strata}）这条线——它用连续特质向量和可累加性区分，把“哪些机器人组合能满足任务需求”形式化为优化问题。在 CBBA 的具体扩展上，文献有几条路线：给 CBBA 加能力约束（最简单，改候选集）；处理需多机协作的耦合任务；把能量/续航纳入代价（energy-aware MRTA，如 2024 年针对长航时动态场景的异构 MRTA 框架\cite{hmrta2024longendurance}，处理电量、充电、任务中继）。本节把这三条路线压缩成对 CBBA 的三个具体改动。

\subsection{理论：异构 CBBA 的三个改动}
\label{sec:hmr-cbba-three}
回顾 CBBA 的两阶段骨架（重述以便不翻回去）：\emph{阶段一（bundle 构建）}每个机器人 $i$ 维护任务包 $b_i$ 和路径 $p_i$，反复把“边际收益最高”的任务加入包中，直到包满或无正收益任务，边际收益为把 $j$ 插入路径最优位置后的得分增量；\emph{阶段二（冲突消解）}每个机器人维护中标价向量 $\mathbf{y}_i$ 与中标者向量 $\mathbf{z}_i$，通过与邻居交换做\emph{带时间戳的最大值共识}，被别人以更高价抢走的任务连同其后任务从包里释放。异构化在这个骨架上做三个改动。\textbf{关键是：阶段二（共识冲突消解）完全不变}——出价是标量，最大值共识对标量的运算与机器人是否异构无关。所有改动都在阶段一（问题定义）。

\paragraph{改动一：能力硬约束——改候选集。}
在 bundle 构建阶段，机器人 $i$ 考虑加入任务 $j$ 之前先检查资格：
\begin{equation}
  i\in\mathcal{A}_j \iff \mathbf{q}_i\succeq\mathbf{y}_j\quad(\text{即 }q_{i,u}\ge y_{j,u},\ \forall u\in\text{所需维度}).
  \label{eq:hmr-eligible}
\end{equation}
只有 $i\in\mathcal{A}_j$ 时 $i$ 才计算 $j$ 的边际收益并可能加入包，否则边际收益直接视为 $-\infty$。这一步是纯本地的、$O(U)$ 的向量比较，几乎不增加计算量，却从根本上杜绝了不可行分配。

\paragraph{改动二：联盟任务——改 bundle 的原子单位。}
有些任务单机能力不足，必须由多机能力拼合，如“搬运 $60\,\mathrm{kg}$ 物体”单台载重 $30\,\mathrm{kg}$ 的四足做不了、需两台组队。这类\textbf{联盟任务（coalition task）}的能力需求 $\mathbf{y}_j$ 在某可累加维上超过任何单机能力。处理有两种思路。\emph{思路 A——任务分解}：把联盟任务拆成多个子任务（“四足$_1$ 抬左端”“四足$_2$ 抬右端”）加时序/同步约束，再用扩展 CBBA 分配子任务；优点是复用单任务 CBBA，缺点是要显式管理子任务间同步（两台必须同时抬）。\emph{思路 B——联盟竞标}：让机器人对“加入某个联盟”出价。任务 $j$ 的合格联盟集合为
\begin{equation}
  \mathcal{C}_j = \Big\{S\subseteq\mathcal{R}:\ \mathrm{agg}\big(S\big)\succeq\mathbf{y}_j\Big\},
  \label{eq:hmr-coalition-set}
\end{equation}
其中 $\mathrm{agg}$ 就是 \cref{eq:hmr-agg} 的聚合（可累加维求和、不可累加维取 $\max$）。机器人竞标“以最小代价满足 $\mathbf{y}_j$ 的联盟中的一席”，更接近 STRATA 视角，但组合复杂度高（联盟数随队伍规模指数增长），实践中要靠贪婪或启发式构造。

\begin{insight}[没有联盟，异构就只是“各干各的”]
若不引入联盟任务、坚持“每个任务必须单机完成”，则所有“重于单机载重”的搬运、所有“需侦察 $+$ 操作两种能力”的复合任务，都会被判为不可分配（候选集为空）——异构系统最有价值的能力（用多机能力拼出单机做不到的事）就被白白浪费。\textbf{联盟任务不是 CBBA 的可选装饰，而是异构协同“$1+1>2$”红利的算法载体。}
\end{insight}

\paragraph{改动三：能量约束——改代价函数。}
续航短的机型（无人机 $\sim 20\,\min$ vs 四足 $\sim 2\,\mathrm{h}$）在分配时必须考虑能耗，否则会出现“无人机接了一串远距离任务，飞到一半没电掉下来”。把能量纳入代价：
\begin{equation}
  c_{ij}^{\mathrm{energy}} = \mathrm{reward}_j - \alpha\cdot\mathrm{energy\_cost}_i(\mathrm{travel}_{ij}) - \beta\cdot\mathbb{1}[\mathrm{energy}_i < \mathrm{threshold}],
  \label{eq:hmr-energy}
\end{equation}
其中 $\mathrm{energy\_cost}_i$ 是机型相关的能耗模型（无人机悬停 $+$ 移动功耗远高于地面车），最后一项是低电量惩罚。更精细的处理（如 2024 年长航时 MRTA 工作\cite{hmrta2024longendurance}）还会显式建模充电站、任务中继。能量约束的本质是给代价函数加一个机型相关、状态相关的修正项——恰好用上能力矩阵里的续航维度。

\subsection{理论-工程桥接：50\% 最优界何时失效}
\label{sec:hmr-cbba-bound}
这是本节最重要的理论洞察，也是异构分配区别于同质 CBBA 的关键代价。CBBA 在 \textbf{DMG（Diminishing Marginal Gain，边际收益递减）}条件下保证不差于最优 $50\%$。DMG 的含义是：一个任务的边际收益不会因为先做了别的任务而增加——形式化地，对任意路径 $p\subseteq p'$ 有 $c_{ij}(p)\ge c_{ij}(p')$。这是一个\textbf{子模性（submodularity）}条件，是顺序贪婪等价于集中式贪婪、从而拿到 $50\%$ 界的数学基础。
\begin{itemize}
  \item \textbf{改动一（能力约束）不破坏 DMG}：它只是缩小候选集，每个机器人在自己的可行任务子集上仍满足 DMG，$50\%$ 界保持。
  \item \textbf{改动三（能量约束）一般不破坏 DMG}：只要能耗是路径长度的、且“已有任务越多、新任务边际收益越小”仍成立（通常成立，电量越用越少、新任务越难加），DMG 保持。
  \item \textbf{改动二（联盟任务）可能破坏 DMG}：这是危险所在。联盟任务引入 agent 间的\emph{正向耦合}——任务 $j$ 对机器人 $i$ 的价值依赖于是否有另一机器人 $k$ 也加入这个联盟。若 $i$ 单独无法完成 $j$（收益为 $0$），但 $i,k$ 一起能完成（收益为正），那么“$k$ 加入”反而\emph{增加}了 $i$ 的边际收益——这正是 DMG（边际收益只减不增）的反面！这种\textbf{超可加（super-additive）效用}违反子模性。
\end{itemize}

\begin{insight}[竞争性分配 vs 合作性分配]
联盟任务把异构协同从“竞争性分配”（大家抢有限任务，多一个竞争者收益降低，子模）变成“合作性分配”（大家凑能力完成任务，多一个合作者收益升高，超可加）。CBBA 的 $50\%$ 界是为竞争性分配设计的，一旦出现超可加的合作收益，这个界就可能失效。\textbf{这不是 CBBA 实现的 bug，而是问题结构的根本变化}——它把我们引向协同搬运（合抬重物，效用超可加）那类“强协同”问题，那里需要的是另一套方法（联合优化、显式联盟形成），而非分布式拍卖。
\end{insight}

\noindent 这给出清晰的工程判据：若异构任务\textbf{都能单机完成}（只是能力筛选 $+$ 能量加权）$\to$ 异构 CBBA 安全、$50\%$ 界保持，放心用分布式拍卖；若存在\textbf{必须多机拼能力}的联盟任务（超可加）$\to$ CBBA 保证失效，要么接受“无保证的启发式”，要么对这部分任务上专门的联盟形成算法。下面给出建立在已有 CBBA 之上的增量异构化核心代码——关键策略是把所有异构逻辑封装进“边际收益计算”和“候选资格判断”，让冲突消解阶段对异构完全无感知（一种关注点分离）。

\begin{codebox}[Python]{异构 CBBA 的 bundle 构建（冲突消解阶段复用，不改）}
import numpy as np

def is_eligible(robot_caps, task_demand, dims):      # 改动一:能力硬约束（所需维度全达标）
    return np.all(robot_caps[dims] >= task_demand[dims])

def marginal_score_hetero(rid, j, path, Q, demands, dims, energy, e_model,
                          alpha=1.0, lowE_pen=1e3, lowE_th=0.2):
    # 改动一:不合格直接 -inf（从竞标剔除），不是大代价软惩罚
    if not is_eligible(Q[rid], demands[j], dims[j]):
        return -np.inf, None
    base_gain, pos = best_insertion_gain(rid, j, path)     # 复用同构 CBBA
    e_cost = e_model(rid, j, path, pos)                    # 改动三:机型相关能耗
    score = base_gain - alpha * e_cost
    if energy[rid] - e_cost < lowE_th:
        score -= lowE_pen                                 # 低电量软回避（不是硬剔除）
    return score, pos

def build_bundle_hetero(rid, tasks, Q, demands, dims, energy, e_model, cap):
    bundle, path = [], []
    while len(bundle) < cap:
        best, best_s, best_p = None, 0.0, None
        for j in tasks:
            if j in bundle: continue
            s, p = marginal_score_hetero(rid, j, path, Q, demands, dims, energy, e_model)
            if s > best_s: best, best_s, best_p = j, s, p
        if best is None: break                            # 无正收益可加任务
        bundle.append(best); path.insert(best_p, best)
    return bundle, path
\end{codebox}

\noindent 三类典型错误：\emph{用大代价软惩罚代替硬剔除}（不合格返回 \texttt{base\_gain - 1e6}——无人合格且“不完成”惩罚更大时仍会被选中，物理不可行分配被派出）；\emph{对联盟任务用单机 CBBA 直接分配}（需求超单机能力时 \texttt{is\_eligible} 对所有单机返回 \texttt{False}、候选集为空、任务静默丢失，须分配后检查未认领任务）；\emph{能量惩罚也用硬剔除}（过度硬化——低电量是“不太好”而非“绝对不可以”，也许这是唯一能做该任务的机器人，宁可让它做完再充电）。\cref{tab:hmr-coalition-methods} 对比三种处理联盟任务的方式。

\begin{table}[H]\centering\small
\caption{三种处理联盟任务的方式（以“搬 $60\,\mathrm{kg}$、单台四足 $30\,\mathrm{kg}$”为例）}
\label{tab:hmr-coalition-methods}
\begin{tabular}{llll}
\toprule
方式 & 复杂度 & DMG / $50\%$ 界 & 适用 \\
\midrule
A 任务分解（拆子任务 $+$ 同步） & 低（复用 CBBA） & 子任务间仍可能超可加，界不保证 & 工程首选 \\
B 联盟竞标（接近 STRATA） & 高（指数，需剪枝） & 显式处理，可上专门保证 & 研究 / 小规模 \\
C 忽略联盟（仅单机分配） & 最低 & 无联盟时保持，有联盟时失败 & 确认无联盟任务 \\
\bottomrule
\end{tabular}
\end{table}

\begin{pitfall}[编程陷阱：能力约束用软惩罚导致不可行分配被选中]
\textbf{错误做法}：把“机器人做不了任务”编码成大代价 $c_{ij}=10^6$ 而非剔除候选。\textbf{现象}：绝大多数实例正常，但在“某任务无合格机器人 $+$ 不完成惩罚更大”的边角实例里，物理上不可能的任务（如让无人机搬 $50\,\mathrm{kg}$）被分配出去、执行端卡死。\textbf{根本原因}：把硬约束（能力，布尔）降级成软约束（代价，连续），连续代价永远可能被“其他更糟选择”反超。\textbf{正确做法}：维护合格候选集 $\mathcal{A}_j$，不合格者边际收益直接 $-\infty$。\textbf{自检}：分配后断言“每个被分配的 $(i,j)$ 都满足 $\mathbf{q}_i\succeq\mathbf{y}_j$”。
\end{pitfall}

\begin{pitfall}[概念误区：“异构 CBBA 也有 50\% 最优界保证”]
\textbf{错误理解}：CBBA 有 $50\%$ 界，加了异构扩展所以异构 CBBA 也有。\textbf{现象}：在含联盟任务的实例上，异构 CBBA 给出的解远差于最优的一半，且无法解释。\textbf{根本原因}：$50\%$ 界依赖 DMG（子模性），联盟任务引入超可加效用（多一个合作者收益升高）直接违反 DMG。\textbf{正确理解}：异构 CBBA 的保证是分情况的——纯能力 $+$ 能量约束保持 $50\%$ 界，联盟任务破坏它。判断标准：任务是“抢”（竞争，子模，有界）还是“凑”（合作，超可加，无界）。
\end{pitfall}

\begin{pitfall}[思维陷阱：把“能量低”和“能力缺失”用同一种约束处理]
\textbf{错误做法}：要么都用硬剔除（电量低也剔除），要么都用软惩罚（能力不足也只罚）。\textbf{现象（都硬）}：电量稍低但仍能完成任务的无人机被剔除、任务无人做（过度保守）；\textbf{现象（都软）}：不会飞的地面车被软性派去航拍（不可行分配泄漏）。\textbf{根本原因}：没区分约束的物理性质——能力缺失是“绝对不可能”（硬），能量低是“有风险但可能可以”（软）。\textbf{正确做法}：能力 $=$ 硬约束（剔除候选）；能量 $=$ 软约束（代价惩罚，留余地）。
\end{pitfall}

\begin{practice}
\begin{enumerate}
  \item （实现 $+$ 实验）在已有 CBBA 上加入改动一与改动三，构造“$3$ UAV（可飞/不可搬重）$+$ $2$ UGV（可搬重/不可飞）”、任务含“侦察”（需飞行）和“搬运”（需载重 $\ge 20\,\mathrm{kg}$）的实例。验证 (a) 侦察只分给 UAV、搬运只分给 UGV；(b) 把某 UAV 初始电量设很低，观察能量惩罚是否让它避开远距离侦察。再做对照：把能力约束改成大代价软惩罚，构造所有 UGV 满载的实例，复现“搬运任务被错误派给 UAV”的 bug。
  \item （推导 $+$ 反例）严格论证纯能力约束的异构 CBBA（改动一）保持 DMG（提示：能力约束只把可行任务集从全集缩小到子集，子集上 DMG 是否保持？）。然后构造一个含联盟任务的最小实例（$2$ 机器人、$1$ 个需两机合作的任务），手动计算边际收益，证明 DMG 被破坏（存在 $p\subseteq p'$ 使 $c_{ij}(p)<c_{ij}(p')$）。
  \item （跨章综合，串联共识与 CBBA）异构 CBBA 的冲突消解沿用最大值共识。假设通信图异构（无人机半径大、地面车小），且无人机可能因续航耗尽中途退出（节点消失）。分析：(a) 持有多个中标权的 UAV 突然退出时，这些任务如何被“撤销”和“重新拍卖”；(b) 最大值共识在节点动态消失时还能保证一致吗、需要什么额外机制（租约/超时）；(c) 哪种通信架构（\cref{sec:hmr-comm}）对“UAV 中途退出”最鲁棒。
\end{enumerate}
\end{practice}

% ════════════════════════════════════════════════════════════════
\section{空地联合通信架构}
\label{sec:hmr-comm}

\paragraph{动机：分配做完了，指令传不到。}
前两节解决了“谁干什么”，但有一个被默认成立、实际却很脆弱的前提：\textbf{分配结果能下发到每个机器人、每个机器人的状态能汇聚回来}。同构队伍里这通常不是问题（大家通信能力一样、距离相近），异构空地队伍里随时可能崩——林区搜救场景：地面车在树林里慢慢搜、和山脚基站直接通信时断时续，无人机在 $100\,\mathrm{m}$ 高空视野开阔、通信通畅，但 $20$ 分钟就要回来换电池。信息到底怎么从基站流到深山里的地面车？谁来当中间人？无人机回去换电池那几分钟系统是不是就断联了？这就是异构通信架构要回答的。它与多机的“三大架构”（集中/分布/去中心）相关但不同：后者关心\emph{决策权}在哪（谁做决定），本节关心\emph{信息流}怎么走（数据怎么传）。异构系统里信息流的设计必须考虑节点异构性——哪类节点适合当中继（视野好、通信半径大）、哪类节点续航是瓶颈、哪类节点失效会导致全网瘫痪。

\subsection{三类架构与单点故障谱}
\label{sec:hmr-comm-three}
把三类架构沿一条“单点故障脆弱性”的谱排列，是理解它们取舍的最佳视角（\cref{fig:hmr-comm}）。

\begin{figure}[ht]
\centering
\begin{tikzpicture}[
  >=Stealth, font=\footnotesize, node distance=6mm,
  n/.style={draw, rounded corners, inner sep=2.5pt, fill=main!8, draw=main!60!black},
  u/.style={draw, rounded corners, inner sep=2.5pt, fill=second!10, draw=second!60!black},
  ed/.style={thick, gray!55!black}]
% --- relay ---
\node[u] (rb) at (0,0) {基站};
\node[n] (ru) at (2.0,0.55) {UAV 中继};
\node[u] (rg) at (4.0,0) {远端 UGV};
\draw[ed,<->] (rb)--(ru); \draw[ed,<->] (ru)--(rg);
\node[align=center,font=\scriptsize] at (2.0,-0.7) {中继型\\（单点故障）};
% --- mesh ---
\begin{scope}[xshift=5.6cm]
\node[n] (m1) at (0,0.7) {U$_1$}; \node[n] (m2) at (1.4,0.7) {U$_2$};
\node[u] (m3) at (0,-0.3) {G$_1$}; \node[u] (m4) at (1.4,-0.3) {G$_2$};
\draw[ed,<->] (m1)--(m2); \draw[ed,<->] (m1)--(m3); \draw[ed,<->] (m2)--(m4);
\draw[ed,<->] (m3)--(m4); \draw[ed,<->] (m1)--(m4); \draw[ed,<->] (m2)--(m3);
\node[align=center,font=\scriptsize] at (0.7,-0.95) {mesh 型\\（无单点故障）};
\end{scope}
% --- leader-follower ---
\begin{scope}[xshift=9.0cm]
\node[u] (l) at (0.8,0.7) {UGV (leader)};
\node[n] (f1) at (0,-0.3) {UAV$_1$}; \node[n] (f2) at (1.6,-0.3) {UAV$_2$};
\draw[ed,->] (l)--(f1); \draw[ed,->] (l)--(f2);
\node[align=center,font=\scriptsize] at (0.8,-0.95) {leader--follower\\（leader 单点，可选举）};
\end{scope}
\end{tikzpicture}
\caption{三类空地通信架构沿单点故障谱排列：中继型把所有依赖押在一个高能力节点（最脆弱）、mesh 型把能力摊平到所有节点（最鲁棒）、leader--follower 居中（leader 失效可选举恢复）。}
\label{fig:hmr-comm}
\end{figure}

\paragraph{架构 A：中继型（relay）——最简单，但最脆弱。}
无人机利用高空视野和大通信半径，充当基站与远端地面车之间的中继。这是空地协同最自然的架构——无人机本来就在高处，让它顺带转发信号几乎零额外成本。\textbf{优势}：实现极简；通信距离大幅延伸（无人机当“移动信号塔”）；地面车不需要强通信能力。\textbf{致命弱点}：\emph{单点故障}——无人机一旦失效（撞树、没电、被干扰），基站和远端地面车立刻完全断联（搜救场景里“无人机回去换电池”那几分钟链路就断了）。

\paragraph{架构 B：mesh 型（ad-hoc mesh）——最鲁棒，但延迟不可控。}
所有节点（无人机、地面车、基站）组成自组织网状网络，消息通过多跳路由传播，这是去中心化思想在通信层的体现。\textbf{优势}：\emph{无单点故障}——任意节点失效消息可绕道，正好补上中继型的致命弱点，最适合需高可靠性的场景（军事、灾难现场）。\textbf{代价}：延迟不确定（多跳路由，跳数随拓扑变化）；网络管理复杂（路由表维护、拓扑动态变化）；带宽利用率低（转发开销）。

\paragraph{架构 C：leader--follower 型——折中，决策集中。}
选一个节点当 leader（通常是计算能力强、续航长的地面车），其他当 follower；leader 负责决策协调，follower 执行并回传感知。\textbf{优势}：结构清晰、协调简单（决策集中）；适合“一个强节点 $+$ 多个弱节点”的异构结构（一台强地面车 $+$ 几架轻无人机），这恰是很多空地队伍的真实形态。\textbf{弱点}：leader 是单点故障，但比中继型温和——follower 间可能还能局部协作，且可设计 leader 选举（leader election）让备份节点顶上。三架构的系统对比见 \cref{tab:hmr-comm-compare}。

\begin{table}[H]\centering\small
\caption{三类空地通信架构对比}
\label{tab:hmr-comm-compare}
\begin{tabular}{llll}
\toprule
维度 & 中继型 & mesh 型 & leader--follower 型 \\
\midrule
实现复杂度 & 最低 & 最高 & 中 \\
单点故障 & 严重（中继挂 $=$ 全断） & 无 & 中（leader 挂可选举恢复） \\
通信延迟 & 低（单/双跳） & 不确定（多跳） & 低（星型） \\
通信距离 & 大（中继延伸） & 大（多跳接力） & 中（受 leader 半径限） \\
适合的异构结构 & 远端弱通信 $+$ 空中中继 & 节点能力相近、要高可靠 & 一强多弱 \\
续航考量 & 中继节点续航是命门 & 分散、无单一命门 & leader 续航是命门 \\
\bottomrule
\end{tabular}
\end{table}

\begin{insight}[架构本质是鲁棒性、延迟、复杂度的三角权衡]
这三种架构本质上是在\textbf{鲁棒性、延迟、复杂度}三者之间做权衡，而权衡的关键变量是异构节点的\emph{能力分布}。中继型把所有鸡蛋放在一个高能力节点的篮子里——能力越集中、效率越高但越脆弱；mesh 型把能力摊平到所有节点——越分散、越鲁棒但越低效。\textbf{没有最好的架构，只有最匹配当前异构结构和任务可靠性需求的架构。}这与“集中 vs 分布 vs 去中心”的权衡是同一哲学在通信层的回响。
\end{insight}

\paragraph{理论-工程桥接：为不同场景选架构。}
\cref{tab:hmr-comm-scenario} 把三类架构落到三个真实场景。注意这些不是死规则——实践中常\emph{混合}使用（一个大型搜救系统可能在局部用 leader--follower，多个局部之间用 mesh 互联），HetSwarm（2025，\cref{sec:hmr-formation} 详述）就是 leader--follower $+$ 局部阻抗链的混合架构。

\begin{table}[H]\centering\small
\caption{为不同场景选通信架构}
\label{tab:hmr-comm-scenario}
\begin{tabular}{lll}
\toprule
场景 & 关键约束 & 推荐架构（理由） \\
\midrule
灾后搜救 & 高可靠性优先（断联 $=$ 救援失败） & mesh（任意节点失效不瘫痪；可靠性 $>$ 延迟） \\
大范围巡检 & 通信距离大、地面车在远端弱通信区 & 中继（无人机当空中中继延伸覆盖） \\
物流配送 & 一个调度中心 $+$ 多个执行单元、结构稳定 & leader--follower（调度中心天然是 leader） \\
\bottomrule
\end{tabular}
\end{table}

\begin{pitfall}[对比性思维：跨类型通信不传私有状态，而是约定公共语义层]
异构系统的跨类型通信\textbf{不是}“把各自的完整状态广播出去让别人解析”，\textbf{而是}“约定一个公共语义层（如世界系目标位姿），各节点只在这一层交换信息”。前者会逼迫每个节点理解所有其他机型的私有状态格式（无人机要懂四足的 $12$ 维关节、四足要懂无人机的姿态四元数），耦合度爆炸；后者让每个节点只需把自己的私有状态投影到公共语义，解耦清晰。这正是前置自测第 $2$ 题“异构状态不能直接共识”的工程对应——通信层必须先有公共语义层。\textbf{自检}：问“加一种全新机型，需要改几个已有节点的代码？”——理想答案是 $0$。
\end{pitfall}

\begin{pitfall}[概念误区：“mesh 网络最鲁棒，所以总是最优选择”]
\textbf{现象}：在“一台强地面车 $+$ $3$ 架轻无人机”的简单巡检里上了 full mesh，结果路由开销吃掉大量带宽、延迟抖动让无人机控制回路不稳，维护这套 mesh 的工程量远超任务本身。\textbf{根本原因}：把“鲁棒性”当成唯一目标，忽略了延迟、复杂度、带宽的代价。\textbf{正确理解}：鲁棒性是有成本的。任务对可靠性要求不高、结构稳定、有天然 leader 时，leader--follower 或中继型更划算——架构选择是多目标权衡，不是单目标最大化。
\end{pitfall}

\begin{pitfall}[思维陷阱：忽略续航差异，让续航短的节点当关键枢纽]
\textbf{错误做法}：选通信能力最强的节点当中继/leader，不管它续航。\textbf{现象}：选了视野最好的无人机当唯一中继，结果它 $20$ 分钟就要回去换电池，每次断联 $5$ 分钟，任务被反复打断。\textbf{根本原因}：只看了“通信能力”这一个异构维度，忽略了“续航”这个维度——异构节点的多个能力维度可能冲突（通信最强的恰好续航最短）。\textbf{正确做法}：选关键枢纽节点时综合考虑通信能力 \emph{且} 续航；若续航短的节点通信最强，要么用 mesh 分散依赖，要么设计“中继接力”（多架无人机轮班当中继）。\textbf{自检}：算“枢纽节点续航 vs 任务时长”——若枢纽续航 $\ll$ 任务时长，必须有接力或冗余机制。
\end{pitfall}

\begin{practice}
\begin{enumerate}
  \item （设计 $+$ 仿真）搭一个“$1$ 强地面车 $+$ $2$ 无人机”的 leader--follower 通信骨架：地面车当 leader 发布共享目标位姿，两架无人机订阅并把模拟“侦察结果”发布到共享话题。注入故障：(a) 模拟 leader 失效，观察系统行为；(b) 实现简单 leader 选举（leader 失联超时后 follower 中 ID 最小者升为新 leader）。对比有无 leader 选举时的恢复能力。
  \item （分析题）对三类架构分别画出“节点失效数 vs 系统连通性”的定性曲线（横轴失效节点数、纵轴仍能通信的节点比例）。哪一类是“阶跃下降”（关键节点失效即崩）、哪一类是“渐进下降”（优雅降级）？这条曲线和 \cref{sec:hmr-axes} 练习 3 分析的“UAV 失效后 $\lambda_2$ 变化”有何联系？
  \item （开放设计题）为一个大型仓库（$4$ 个区域，每区 $1$ AGV $+$ $2$ 盘点无人机、共享一个调度中心）设计\emph{混合架构}：区域内用什么、区域间用什么、调度中心如何接入？某区域 AGV 故障时如何让该区无人机继续工作（被邻区接管？直连调度中心？）？画出拓扑并标注每条链路的架构类型。
\end{enumerate}
\end{practice}

% ════════════════════════════════════════════════════════════════
\section{异构编队：leader--follower 与阻抗引导}
\label{sec:hmr-formation}

\paragraph{动机：同一条参考轨迹，喂不进异构编队。}
分布式 MPC 编队很美：$N$ 个机器人用 ADMM 把各自的单体 MPC 耦合起来，邻居之间维持期望相对位置 $p_i-p_j=d_{ij}$，整个队形像刚体一样移动。这套方法的前提是所有机器人共享同一动力学模型模板（如都是二阶积分 $\ddot p=u$，只是质量/推力上限不同），一致性约束 $p_i-p_j=d_{ij}$ 直接写在它们共有的位置状态上，干净利落。现在把这个前提打破，考虑一个编队里有：一架四旋翼（欠驱动，水平加速靠倾斜，不能瞬间侧移）、一台差速车（非完整，瞬时速度锁在车头朝向，不能横移）、一个全向平台（可任意方向移动）。让它们保持三角形队形向前移动——编队需要向左平移一小步：全向平台直接向左走；四旋翼得先向左倾斜再产生向左加速度（位置能到、有动态延迟）；差速车\emph{根本不能直接向左走}（非完整约束禁止横移，必须先转向再前进，走曲线才能到达左边目标）。如果硬套同构方法把一致性约束写在状态空间，会发现三种机器人状态维度都不同（四旋翼 $12$、差速车 $3$、全向 $4$），“相对状态相等”\emph{根本无法定义}（没法让 $12$ 维向量减 $3$ 维向量）；即便强行只比较位置子空间，约束 $p_i-p_j=d_{ij}$ 对差速车也\emph{瞬时不可满足}。这就是异构编队的根本困难：\textbf{异构动力学的状态空间不可比，可达集形状不同，状态空间一致性失效。}

\paragraph{反面：如果硬在状态空间做一致性会怎样。}
强行在位置子空间写一致性约束 $p_i-p_j=d_{ij}$ 喂给每个 agent 的局部 MPC，会发生三件事：（1）\textbf{差速车的 MPC 不可行（infeasible）}——某些时刻“达到期望相对位置”要求差速车瞬间横移、违反非完整约束、求解器报 infeasible、控制断流；（2）\textbf{四旋翼的姿态被忽略}——位置一致性完全不管姿态，但四旋翼位置-姿态强耦合，只约束位置可能逼出超出姿态控制能力的剧烈机动；（3）\textbf{整个编队被“最弱者”拖累且不自知}——即便勉强可行，编队的可达运动被差速车非完整约束限制，但状态空间写法看不到这一点，规划出的队形机动差速车跟不上、队形持续散开。根本原因：\textbf{状态空间一致性隐含假设“所有 agent 能瞬间达到任意相对状态”，这只对全向、全驱动系统成立。}异构编队里有非完整、欠驱动成员，这个假设崩了。正确出路是把一致性\emph{上移一层}——不在状态空间约束，而在\textbf{输出空间（output space）}约束。

\paragraph{历史：从输出调节到异构编队。}
把多机协调建立在“输出”而非“状态”上，思想源头是控制理论里的\textbf{输出调节（output regulation）}和\textbf{输出一致性（output consensus）}。核心观念是：异构系统可以有不同的内部状态和动力学，但只要它们的\emph{输出}（我们关心的可观测量，如机体在世界系的位置）能达成一致，协调目标就实现了。每个 agent 用自己的（异构的）动力学和控制器去驱动自己的输出，输出层面才做耦合。在机器人编队里这条线发展出“虚拟结构（virtual structure）”和“leader--follower 输出跟踪”等方法。近年的 HetSwarm（2025）\cite{hetswarm2025} 给出了一个特别优雅的工程实例：它不试图让异构成员保持刚性几何，而是用\textbf{阻抗（impedance）}把它们柔性地连起来——leader 在输出空间规划，follower 通过阻抗链跟随，且 follower 能根据自己的可达集对障碍做出与 leader 不同的响应。

\subsection{理论：公共输出空间一致性}
\label{sec:hmr-formation-output}
异构编队的设计分三步走，核心是“找到公共输出空间”。

\paragraph{第一步：定义公共输出空间。}
为每个 agent $i$ 定义输出映射 $h_i$，把它的（异构）状态 $x_i$ 投影到一个\emph{所有 agent 共享的、同维同义的}输出空间：
\begin{equation}
  y_i = h_i(x_i) \in \mathbb{R}^p,\qquad \forall i.
  \label{eq:hmr-output-map}
\end{equation}
关键是 $h_i$ 因 agent 而异（异构），但输出 $y_i$ 同维同义（公共）。最常用的公共输出是\textbf{机体在世界系下的位置}（2D 地面编队取 $(x,y)$，3D 空地编队取 $(x,y,z)$）：四旋翼 $h_{\mathrm{uav}}(x)=$ 从 $12$ 维状态取位置分量；差速车 $h_{\mathrm{ugv}}(x)=(x,y)$；全向平台 $h_{\mathrm{omni}}(x)=$ 位置分量。现在所有 agent 的输出都是世界系位置，\emph{同维、同义、可比较、可做共识}——前置自测第 $2$ 题的答案在这里落地：异构状态不能直接共识，但投影到公共输出后可以。

\paragraph{第二步：在输出空间写一致性。}
编队约束写在输出空间：
\begin{equation}
  y_i - y_j = d_{ij}\qquad(\text{期望的输出相对几何，即队形}).
  \label{eq:hmr-output-consensus}
\end{equation}
这个约束现在是良定义的（同维向量相减），且对所有 agent 一视同仁，描述的是“我们想要的队形几何”，不涉及任何 agent 的内部动力学。

\paragraph{第三步：每个 agent 用各自动力学实现输出目标。}
一致性给出每个 agent 的\emph{输出目标} $y_i^\ast$（为维持队形它应到达的世界系位置），然后\textbf{每个 agent 用自己的（异构的）局部控制器}去实现：四旋翼用微分平坦 $+$ $\SE(3)$ 控制生成可行轨迹到达 $y_i^\ast$；差速车用考虑非完整约束的轨迹跟踪（pure pursuit / 非完整 MPC）到达 $y_i^\ast$；全向平台直接位置控制。

\begin{insight}[在公共输出空间协调，在私有状态空间执行]
异构编队的核心设计原则可以浓缩成一句话——\textbf{“在公共输出空间协调，在私有状态空间执行”}。协调层只关心“每个 agent 的输出应该在哪”（队形几何，所有 agent 同语言），执行层让每个 agent 用自己的动力学去实现（各说各话）。这是关注点分离在异构编队上的体现，也是它区别于同构编队的根本之处：同构编队协调和执行都在同一个状态空间（因为同构），异构编队必须把这两层劈开。这个“协调/执行分离”会在 \cref{sec:hmr-marl} 的异构 MARL（共享的高层 $+$ 私有的低层动作）和 \cref{sec:hmr-bound} 的分层架构里再次出现——它是异构系统设计的元原则。
\end{insight}

\subsection{改进 APF 与阻抗引导}
\label{sec:hmr-formation-impedance}
HetSwarm 给出的具体实现是 leader--follower $+$ 改进人工势场（APF）。\textbf{Leader（无人机）}用 APF 在输出空间规划：目标产生引力、障碍产生斥力，合力引导 leader 实时调整轨迹（APF 反应式、计算轻、天然处理动态障碍）。\textbf{Follower（地面车）}通过\textbf{阻抗链（impedance link）}跟随 leader：把 leader--follower 之间的期望相对位置当成虚拟弹簧-阻尼系统，follower 受到一个把它拉向期望队形位置的“阻抗力”。阻抗模型把跟随写成二阶系统\cite{hogan1985impedance}：
\begin{equation}
  M_d(\ddot y_i - \ddot y_i^\ast) + B_d(\dot y_i - \dot y_i^\ast) + K_d(y_i - y_i^\ast) = F_{\mathrm{ext}},
  \label{eq:hmr-impedance}
\end{equation}
其中 $M_d,B_d,K_d$ 是虚拟惯量/阻尼/刚度，$y_i^\ast$ 是 leader 给出的期望输出，$F_{\mathrm{ext}}$ 是外部交互力（如局部避障）。这个柔性连接比刚性约束 $y_i-y_j=d_{ij}$ 宽容得多——它允许 follower 在遇到障碍时\emph{暂时偏离}期望队形（弹簧被拉长），过后再回弹。

\paragraph{阻抗引导处理异构可达集差异。}
HetSwarm 最妙的一点正好踩中异构编队的痛点。考虑一个矮障碍（地面上的一块石头）：对无人机（leader）它在高空飞，\emph{这个矮障碍根本不在它的飞行高度}——可以无视；对地面车（follower）这个矮障碍\emph{挡在路上}、必须绕开。如果用刚性队形约束，地面车绕障碍时会破坏队形、进而通过刚性耦合干扰无人机。HetSwarm 的做法是：让地面车对这类矮障碍建立\textbf{临时阻抗链（temporal impedance link）}——障碍施加临时斥力，地面车局部绕开（阻抗弹簧被拉伸），而无人机轨迹完全不受影响（障碍不在它的可达集相关区域）。

\begin{insight}[对比性思维：每个 agent 只对自己可达集相关的障碍反应]
异构编队的避障\textbf{不是}“整个编队一起绕开所有障碍”，\textbf{而是}“每个 agent 只对自己可达集相关的障碍做出反应”。无人机不该为地面的石头绕路（它能飞过去），地面车不该为高处的树枝绕路（它在树下）。刚性编队强迫所有成员对所有障碍统一反应，浪费且低效；阻抗引导让每个成员根据自己的可达集选择性地反应，柔性连接吸收了不一致。\textbf{这把异构性从“必须统一妥协的负担”变成了“各扬所长的红利”}——无人机的飞行能力让它免于地面绕障，正是异构的价值所在。
\end{insight}

\noindent 下面给一个 2D 的最小实现：无人机（leader，全向近似）用 APF 规划、差速车（follower，非完整）用阻抗 $+$ 非完整跟踪。分层逻辑是：输出空间（世界系位置）是 leader 和 follower 唯一能“对话”的公共语言（它们状态空间不可比，但位置可比）；APF 给 leader（反应式、轻量），阻抗给 follower（柔性连接允许临时偏离队形去绕障），非完整跟踪是执行层（阻抗给出“期望位置/速度”这个输出空间目标，差速车再用自己的非完整控制器去实现——这就是“协调/执行分离”）。

\begin{codebox}[Python]{异构 leader--follower 编队（APF $+$ 阻抗 $+$ 非完整跟踪）}
import numpy as np

def apf_velocity(pos, goal, obstacles, k_att=1.0, k_rep=2.0, rep_radius=2.0):
    """Leader 人工势场速度（输出空间=世界系位置）:目标引力 + 障碍斥力."""
    to_goal = goal - pos; dist = np.linalg.norm(to_goal) + 1e-9
    f_att = k_att * to_goal / dist * min(dist, 1.0)       # 远处饱和避免过冲
    f_rep = np.zeros(2)
    for obs in obstacles:
        diff = pos - obs; d = np.linalg.norm(diff) + 1e-9
        if d < rep_radius:                                # 斥力必须有有限影响半径
            f_rep += k_rep * (1.0/d - 1.0/rep_radius) / (d**2) * (diff / d)
    return f_att + f_rep

def impedance_target(fpos, fvel, des_pos, des_vel, local_obs,
                     Md=1.0, Bd=4.0, Kd=6.0, k_obs=3.0, obs_radius=1.0):
    """Follower 阻抗:拉向期望队形位置（虚拟弹簧-阻尼）+ 对【自己可达集相关的】
       局部障碍建临时阻抗.返回期望加速度（交给非完整跟踪器实现——执行层）."""
    f_imp = Kd * (des_pos - fpos) + Bd * (des_vel - fvel)
    f_obs = np.zeros(2)
    for obs in local_obs:                                 # 仅对地面障碍反应（无人机会忽略的矮障碍）
        diff = fpos - obs; d = np.linalg.norm(diff) + 1e-9
        if d < obs_radius:
            f_obs += k_obs * (1.0/d - 1.0/obs_radius) / (d**2) * (diff / d)
    return (f_imp + f_obs) / Md

def unicycle_track(state, acc_des, dt, k_v=1.0, k_w=2.0):
    """执行层:差速车非完整跟踪器，把输出空间期望加速度转成 (线加速, 角速度)."""
    x, y, theta, v = state
    des_dir = np.arctan2(acc_des[1], acc_des[0])
    herr = np.arctan2(np.sin(des_dir - theta), np.cos(des_dir - theta))
    a = k_v * np.linalg.norm(acc_des) * np.cos(herr)      # 只沿车头方向加速
    omega = k_w * herr
    v_new = v + a * dt; theta_new = theta + omega * dt
    return np.array([x + v_new*np.cos(theta_new)*dt,      # 非完整运动学
                     y + v_new*np.sin(theta_new)*dt, theta_new, v_new])
\end{codebox}

\noindent 三类典型错误：\emph{在状态空间写刚性一致性喂给非完整 follower}（直接要求差速车瞬间到达 leader $+$ offset 的位置当速度指令——当 target 在车侧方时要求横向速度、非完整下无法实现、跟踪器饱和或发散）；\emph{让 follower 对所有障碍反应}（包括无人机才需避的高障碍——地面车为高处树枝绕路，但它在树下、本不该受影响）；\emph{APF 斥力没有影响半径}（远处障碍仍施加微弱斥力、累积让 leader 轨迹持续偏移、甚至陷入局部极小）。\cref{tab:hmr-rigid-soft} 对比刚性约束与阻抗柔性连接。

\begin{table}[H]\centering\small
\caption{异构编队：刚性约束 vs 阻抗（柔性）连接}
\label{tab:hmr-rigid-soft}
\begin{tabular}{lll}
\toprule
维度 & 刚性一致性 $y_i-y_j=d_{ij}$ & 阻抗柔性连接 \\
\midrule
非完整 agent & 不可行（要求瞬时横移） & 可行（执行层用非完整跟踪） \\
遇障碍 & 整队妥协或求解失败 & 局部偏离（弹簧拉伸）后回弹 \\
异构可达集 & 被最弱者拖累 & 各按自己可达集反应 \\
队形精度 & 高（若可行） & 略松（弹簧形变），但鲁棒 \\
计算 & 需联合优化（ADMM） & 轻量（局部反应式） \\
\bottomrule
\end{tabular}
\end{table}

\begin{pitfall}[编程陷阱：把输出空间目标当状态空间指令直接下发给异构 agent]
\textbf{错误做法}：一致性算出 follower 的期望世界系位置，直接当速度/位置指令喂给非完整车。\textbf{现象}：当期望位置在车侧方时差速车无法直接横移、跟踪器饱和或发散、车原地打转。\textbf{根本原因}：跳过了“执行层”——输出空间目标必须经由 agent 自己的（异构）控制器转成它动力学可行的指令，不能绕过动力学直接当指令。\textbf{正确做法}：协调层只给“期望输出”（世界系位置/速度），每个 agent 用自己的控制器（非完整跟踪/微分平坦）去实现。\textbf{自检}：检查下发给 agent 的指令是否是它动力学的合法输入（差速车应是 $(v,\omega)$ 或可达加速度，不是任意方向的位置跳变）。
\end{pitfall}

\begin{pitfall}[概念误区：“异构编队就是给不同机器人设不同增益的同构编队”]
\textbf{错误理解}：同构编队加个增益调度、对不同机型调不同 PID 参数就成异构编队了。\textbf{现象}：增益怎么调，非完整车在侧向机动时都跟不上、队形持续散开。\textbf{根本原因}：增益调度只能处理参数异构（同模型不同参数），处理不了结构异构（不同模型、不同约束、不同可达集）——非完整 vs 全向是结构差异，不是参数差异。\textbf{正确理解}：异构编队的关键不是调参，而是架构改变——把一致性从状态空间上移到输出空间，劈成“协调层 $+$ 各自执行层”。
\end{pitfall}

\begin{pitfall}[思维陷阱：追求刚性队形精度，牺牲对异构可达集的适应]
\textbf{错误做法}：为了队形“好看”，用刚性约束逼所有 agent 严格保持几何。\textbf{现象}：地面车为维持与无人机的精确相对位置被迫做非完整车做不好的机动，或为无人机能飞过的障碍绕远路，整体效率暴跌。\textbf{根本原因}：把“队形精度”当成目标本身，忘了队形是为任务服务的手段——异构成员可达集不同，强求刚性精度等于强求“削足适履”。\textbf{正确理解}：异构编队应优先柔性（阻抗）——允许队形在遇障/机动时临时形变，换取每个成员都能扬长避短。队形是约束，不是目标。
\end{pitfall}

\begin{practice}
\begin{enumerate}
  \item （实现 $+$ 实验）搭一个 2D 异构编队：$1$ leader（APF，全向近似）$+$ $2$ follower（$1$ 差速车用 \texttt{unicycle\_track}、$1$ 全向平台直接位置控制）。让编队穿过有“矮障碍”（只挡地面车）的走廊。验证 (a) 全向平台和差速车都绕开矮障碍、但 leader 轨迹不受影响；(b) 对比用刚性约束时差速车能否跟上。把阻抗刚度 $K_d$ 从小调到大，观察“柔性→刚性”过渡时差速车跟踪质量和队形保持的此消彼长。
  \item （推导题）证明：对全向、全驱动 agent（$\ddot p=u$，输出 $y=p$），输出空间一致性 $y_i-y_j=d_{ij}$ 与状态空间一致性 $p_i-p_j=d_{ij}$ 等价（这说明同构编队是异构编队的特例）。然后说明对非完整 agent 为什么这个等价性破裂（提示：输出映射 $h$ 是否可逆、可达集是否满秩）。
  \item （跨章综合，串联共识/ADMM/输出一致性）共识、ADMM 编队、输出空间一致性本质上都是“让一群 agent 在某个量上达成协调”。用一张表对比：(a) 各自“协调的对象”是什么（标量平均 / 状态相对位置 / 输出相对几何）；(b) 各自对 agent 同构性的要求；(c) 当 agent 异构时共识和 ADMM 各需做什么改造。把分析与“协调/执行分离”元原则联系起来。
\end{enumerate}
\end{practice}

% ════════════════════════════════════════════════════════════════
\section{跨模态信息融合}
\label{sec:hmr-perception}

\paragraph{动机：各自定位很准，拼在一起全错位。}
异构队伍要协同，前提是它们对世界有\emph{一致的认识}——无人机看到的目标和地面车看到的同一个目标，必须能对应到同一个世界坐标。一个看似显然的方案：让无人机和地面车各自跑 SLAM（各自定位 $+$ 建图），把各自估计的位姿广播出去，按位姿拼地图。在\textbf{同构}队伍里（两台一样的地面车）这常能凑合工作——视角相近、传感器相同、有大量共视特征、回环检测能对齐两个坐标系。但在\textbf{异构空地}队伍里会惨败：无人机在 $80\,\mathrm{m}$ 高空下视，看到屋顶、道路、树冠——一片俯视纹理；地面车在地面平视，看到墙面、门窗、家具、行人腿——一片侧视结构。\textbf{这两套观测之间几乎没有任何共视特征}：屋顶和墙面是同一栋楼的不同面，道路俯视纹理和地面车看到的路面侧视完全不同。即便用最强的深度学习描述子也匹配不上——同一物理点在鸟瞰和平视下的外观差异远超描述子的不变性范围。没有共视特征 $\to$ 无法回环 $\to$ 两个独立 SLAM 的坐标系无法对齐 $\to$ 各自位姿在各自的、漂移的、互不相干的坐标系里 $\to$ 广播出来的位姿拼不到一起。结果就是“各自都很准，拼起来差两米”。

\subsection{两道鸿沟：视角鸿沟与模态鸿沟}
\label{sec:hmr-perception-gaps}
\textbf{视角鸿沟（viewpoint gap）}：同一物体/场景在鸟瞰和平视下的外观差异巨大，导致跨视角的特征匹配失效。这是几何层面的鸿沟——观测的是同一个三维世界，但投影视角差了近 $90^\circ$，二维外观几乎无重叠。\textbf{模态鸿沟（modality gap）}：不同传感器（相机 vs 激光雷达 vs 热成像）产生的数据类型根本不同（图像 vs 点云 vs 温度场），无法直接比对。这是数据层面的鸿沟——连数据结构都不一样，谈何匹配。两道鸿沟可单独出现，也可叠加（无人机相机 vs 地面车激光，视角 $+$ 模态双鸿沟）。它们的存在使得“位姿广播 $+$ 地图拼接”这个对同构系统有效的方案，对异构系统\textbf{结构性地失效}。

\begin{insight}[异构融合的真正难点在“鸿沟”，不在“融合”]
同构多机的信息融合和异构多机的信息融合，\textbf{难度不是量的差异而是质的差异}。同构融合的核心挑战是“数据关联 $+$ 一致性估计”（找到共视、融合估计），是已被充分研究的问题；异构融合在数据关联之前先要跨越“鸿沟”——当根本没有共视特征、连数据类型都不同时，传统数据关联无从下手。\textbf{异构融合的真正难点在“鸿沟”，不在“融合”}——一旦跨越鸿沟建立了对应，后面的融合反而是相对成熟的部分。这就是为什么本节是本章最硬核的一节：它要解决的是一类同构系统里根本不存在的问题。
\end{insight}

\paragraph{历史：从协同 SLAM 到跨模态定位。}
多机感知融合的发展可看作“逐步跨越鸿沟”的历史。\emph{同构协同 SLAM}（2010s）：多个相同/相近机器人共享特征、联合优化位姿图，前提是大量共视（无鸿沟或弱鸿沟），代表如多机 ORB-SLAM、CoSLAM。\emph{$2.5\mathrm{D}$ 地图配准}（2010s 中后期）：意识到鸟瞰和平视无法直接特征匹配，转而在\textbf{几何中间表示}上对齐——把视觉重建的 3D 点云和激光的 3D 点云都投影成 $2.5\mathrm{D}$ 高程图（elevation map），用 ICP 对齐高程图，绕过外观鸿沟、改在几何上对齐\cite{aghs25dmap2017}。\emph{深度学习跨视角描述子}（2020s）：用学习把鸟瞰和平视图像编码到\textbf{共享嵌入空间}，让同一地点的两类描述子在嵌入空间里接近，直接攻克视角鸿沟，代表如跨视角地理定位、空地 place recognition。\emph{半分布式相对定位}（2025）：最新思路是\textbf{解耦}——不强求所有机器人在一个全局坐标系里，而是估计两两之间的相对位姿，且把相对定位与各自的状态估计解耦\cite{crossmodal2025semidist}。这条线的主线是：从“消除鸿沟”（找共视）走向“绕过鸿沟”（几何中间表示）再走向“跨越鸿沟”（学习共享嵌入）和“回避全局对齐”（只求相对位姿）。

\subsection{三条技术路线}
\label{sec:hmr-perception-routes}

\paragraph{路线一：$2.5\mathrm{D}$ 高程图配准——在几何上绕过外观鸿沟。}
核心思想：既然鸟瞰和平视的\emph{外观}无法匹配，那就退到它们共有的\emph{几何}上对齐——三维世界的几何结构（地形高度、建筑轮廓）是视角无关的。流程：（1）各自生成 3D 点云（无人机下视相机做多视图几何、地面车激光累积）；（2）投影成 $2.5\mathrm{D}$ 高程图（把 3D 点云投影到水平网格、每格存高度值，“$2.5\mathrm{D}$” $=$ 2D 网格 $+$ 每格一个高度，介于 2D 和完整 3D 之间，高程图是视角无关的）；（3）ICP 对齐（把地面车高程图配准到无人机高程图，求两坐标系相对变换 $\mathbf{T}$）；（4）融合定位（有了 $\mathbf{T}$，两机器人位姿被锚定到同一坐标系，可互相提供定位约束）。\textbf{为什么有效}：屋顶纹理和墙面纹理无法匹配，但“这片区域有一栋 $10\,\mathrm{m}$ 高的楼”这个几何事实无人机和地面车都能观测到——几何是它们的公共语言。\textbf{局限}：需要场景有足够几何特征（高度变化）；平坦无特征区域（空旷停车场）高程图退化、ICP 无法收敛；点云生成本身有误差。

\begin{insight}[多视角理解：高程图就是机器人版的等高线地形图]
$2.5\mathrm{D}$ 高程图配准的思路和读“等高线地形图”高度相似——无论你站在山顶俯瞰还是站在山脚仰望，那张等高线图描述的“哪里高、哪里低”是同一个客观事实，与你的观察视角无关。高程图就是机器人版的等高线图：它剥离了“从哪个角度看”（外观随视角剧变），只保留“这片地形的高度分布”（几何随视角不变）。但类比的边界在于：等高线图是人预先测绘好的静态产物，而机器人的高程图是各自实时从点云投影出来的、带噪声和缺失（未观测区为 NaN）的估计——所以还需要 ICP 去对齐两份不完美的高程图，而不是像查同一张地图那样直接共享。不要把“高程图是公共语言”误解成“两个机器人查的是同一份地图”。
\end{insight}

\begin{insight}[理论-工程桥接：异构融合的艺术是“选对中间表示”]
$2.5\mathrm{D}$ 配准用 ICP，而 ICP 的本质是在两个点集间反复“找最近点对应 $+$ 求最优刚体变换”。这里复用得恰到好处——异构感知融合\emph{没有发明新的几何工具}，而是把成熟的 ICP 用在了“高程图”这个精心选择的、视角无关的中间表示上。\textbf{异构融合的工程艺术，很大程度上是“选对中间表示”的艺术}——选一个让异构观测变得可比的公共表示（高程图就是这样的表示），后面的对齐就能复用成熟工具。
\end{insight}

\paragraph{路线二：深度学习跨视角描述子——直接跨越视角鸿沟。}
核心思想：用神经网络学一个编码器，把鸟瞰和平视图像都映射到\textbf{共享嵌入空间}，使同一地点的两类描述子在嵌入空间里距离近、不同地点的远。流程：（1）训练跨视角编码器（用配对数据——同一地点的鸟瞰 $+$ 平视图像对，损失如对比损失/三元组损失拉近正样本、推远负样本，网络学会忽略视角差异、抓住地点本质）；（2）描述子检索（地面车提取当前平视图像的全局描述子，在无人机鸟瞰描述子数据库里检索最近邻，找到“地面车现在在无人机地图的哪个位置”）；（3）位姿求解（检索到对应后用几何方法 PnP/配准求精确相对位姿）。近期工作（2025 半分布式相对定位\cite{crossmodal2025semidist}）进一步把它做成增量式：UGV 和 UAV 各自提取基于深度学习的关键点 $+$ 全局描述子，UGV 用局部 BA（LiDAR $+$ 相机 $+$ IMU）快速获得精确相对位姿，增量式回环维护和检索关键帧；关键创新是\textbf{把相对定位与所有 agent 的状态估计解耦}。\textbf{局限}：需训练数据（配对的跨视角图像、获取成本高）；泛化性受训练分布限制；对模态鸿沟（相机 vs 激光）不直接适用，需额外设计。

\paragraph{路线三：相对定位解耦——回避全局对齐难题。}
核心思想：很多协同任务其实\emph{不需要}全局一致坐标系，只需要知道机器人两两之间的相对位姿（如“无人机引导地面车到目标”，只要无人机知道目标相对地面车的位姿即可）。把问题从“全局定位”降级为“相对定位”有两个好处：（1）\textbf{避开全局对齐的病态}（全局对齐要求消除所有 agent 累积漂移并对齐到公共原点，异构鸿沟下极难；相对定位只需局部、瞬时的相对关系，宽容得多）；（2）\textbf{解耦状态估计}（每个 agent 独立做自己的 SLAM，相对定位作为独立、解耦的模块，输入各自 SLAM 输出和跨模态观测、输出相对位姿——这就是 2025 半分布式工作的“semi-distributed”：估计是分布式的，但相对定位有一个解耦的协调环节）。\textbf{为什么有效}：它承认并回避了异构全局对齐的根本困难，把目标降到“够用就好”——很多协同任务确实只需相对关系，这是一种务实的工程智慧。三条路线对比见 \cref{tab:hmr-fusion-routes}。

\begin{table}[H]\centering\small
\caption{跨模态融合三条路线对比}
\label{tab:hmr-fusion-routes}
\begin{tabular}{lllll}
\toprule
路线 & 攻克的鸿沟 & 核心工具 & 需训练 & 局限 \\
\midrule
$2.5\mathrm{D}$ 高程图配准 & 视角（绕过到几何） & ICP $+$ 高程图 & 否 & 需几何特征；平坦区退化 \\
深度跨视角描述子 & 视角（直接跨越） & 神经网络嵌入 $+$ 检索 & 是 & 需配对数据；泛化受限 \\
相对定位解耦 & 回避全局对齐 & 局部 BA $+$ 解耦协调 & 部分 & 只给相对位姿、无全局系 \\
\bottomrule
\end{tabular}
\end{table}

\begin{insight}[对比性思维：不一定要全局一致地图，往往只需协作所需的相对一致]
异构感知融合的目标\textbf{不一定是}“建立一个所有机器人共享的、全局一致的世界地图”，\textbf{而往往只是}“让需要协作的机器人对它们协作所需的那部分世界达成一致”。前者是 SLAM 式的全局野心（在异构鸿沟下极难、常常过度），后者是任务驱动的务实（只对齐协作必需的相对关系）。想清楚“我的任务到底需要全局一致，还是只需相对一致”，是选对融合路线的第一步——很多时候路线三就够了，根本不必去啃路线一、二的全局对齐硬骨头。
\end{insight}

\noindent 下面给出 $2.5\mathrm{D}$ 高程图配准的核心骨架。为什么不直接 ICP 原始点云、而要先投影成高程图？直接 ICP 两个原始点云会失败：密度天差地别（无人机视觉点云稀疏不均、地面激光稠密——ICP 的最近点对应被密度差污染）；覆盖不同（无人机看屋顶、地面车看墙面下部，大量点无对应、ICP 被无对应点带偏）。高程图把两者归约到“每个水平网格的高度”——密度被网格化抹平，只比较共同覆盖的水平区域的高度分布，让 ICP 在干净、可比的表示上工作。

\begin{codebox}[Python]{$2.5$D 高程图配准（点云投影 $+$ 2D ICP，实际 3D 用 Open3D）}
import numpy as np

def pointcloud_to_elevation(points, grid_res=0.5, bounds=None):
    """3D 点云投影成 2.5D 高程图:每个水平网格存最高点高度.NaN=未观测."""
    xy_min, xy_max = (points[:, :2].min(0), points[:, :2].max(0)) if bounds is None else bounds
    size = np.ceil((xy_max - xy_min) / grid_res).astype(int) + 1
    elev = np.full((size[1], size[0]), np.nan)
    idx = ((points[:, :2] - xy_min) / grid_res).astype(int)
    for (cx, cy), z in zip(idx, points[:, 2]):
        if np.isnan(elev[cy, cx]) or z > elev[cy, cx]:
            elev[cy, cx] = z                              # 取最高（屋顶/树冠）
    return elev, xy_min

def icp_2d(src, dst, max_iter=50, tol=1e-5):
    """2D ICP:把 src 配准到 dst，返回 R, t（两坐标系相对变换）."""
    from scipy.spatial import cKDTree
    R, t = np.eye(2), np.zeros(2); src_cur = src.copy(); prev = np.inf
    for _ in range(max_iter):
        dists, idx = cKDTree(dst).query(src_cur)         # 最近点对应
        matched = dst[idx]
        sc, dc = src_cur - src_cur.mean(0), matched - matched.mean(0)
        U, _, Vt = np.linalg.svd(sc.T @ dc)              # Kabsch 求最优刚体变换
        Rs = Vt.T @ U.T
        if np.linalg.det(Rs) < 0: Vt[-1] *= -1; Rs = Vt.T @ U.T   # 防反射
        ts = matched.mean(0) - Rs @ src_cur.mean(0)
        src_cur = (Rs @ src_cur.T).T + ts
        R, t = Rs @ R, Rs @ t + ts
        err = dists.mean()
        if abs(prev - err) < tol: break
        prev = err
    return R, t                                          # 地面车系 -> 无人机系
\end{codebox}

\noindent 三类典型错误：\emph{直接广播位姿、按各自 SLAM 坐标拼接}（两个独立 SLAM 的原点、朝向、尺度都不同且各自漂移，直接拼错位严重——“各自准、拼起来差两米”的根源）；\emph{用手工图像描述子跨视角匹配}（ORB/SIFT 的视角不变性远不足以跨越 $90^\circ$ 视角差，屋顶和墙面描述子根本对不上、匹配全是噪声）；\emph{ICP 直接配准原始点云}（密度差和覆盖差让最近点对应大量错误、ICP 收敛到错误解或发散）。\cref{tab:hmr-fusion-choice} 给出按场景条件选融合路线的决策。

\begin{table}[H]\centering\small
\caption{按场景条件选跨模态融合路线}
\label{tab:hmr-fusion-choice}
\begin{tabular}{ll}
\toprule
条件 & 推荐路线 \\
\midrule
只需相对位姿 & 路线三：相对定位解耦（回避全局对齐） \\
需全局系 $+$ 有几何 $+$ 跨模态 & 路线一：$2.5\mathrm{D}$ 高程图配准（几何中间表示，跨模态友好） \\
需全局系 $+$ 有训练数据 $+$ 同模态 & 路线二：深度跨视角描述子（直接跨越视角鸿沟） \\
需全局系 $+$ 有几何（通用兜底） & 路线一：$2.5\mathrm{D}$ 配准（最通用、无需训练） \\
啥都没有 & 加人工标志（fiducial markers）或 GPS 锚点降低难度 \\
\bottomrule
\end{tabular}
\end{table}

\begin{pitfall}[编程陷阱：用手工图像描述子做跨视角匹配]
\textbf{错误做法}：对鸟瞰图和平视图各提 ORB/SIFT、用暴力匹配求对应。\textbf{现象}：匹配点几乎全是错的、求出的相对位姿完全错误、融合后地图错位更严重。\textbf{根本原因}：手工描述子的视角不变性是为“小视角变化”设计的，远不足以跨越鸟瞰 vs 平视的近 $90^\circ$ 视角鸿沟——同一物理点的鸟瞰和平视外观无重叠。\textbf{正确做法}：用视角无关的几何中间表示（$2.5\mathrm{D}$ 高程图 $+$ ICP，路线一），或专门训练的跨视角学习描述子（路线二）。\textbf{自检}：可视化匹配点对，若匹配线杂乱无章（不符合任何刚体变换）即是跨视角失效。
\end{pitfall}

\begin{pitfall}[概念误区：“异构融合就是把各自的位姿广播出来拼一拼”]
\textbf{错误理解}：每个机器人都能自定位，广播位姿后按位姿拼地图即可。\textbf{现象}：各自定位精度都很高（cm 级），但拼接后两图错位米级、无法协作。\textbf{根本原因}：忽略了“坐标系对齐”——各自 SLAM 在各自的、独立漂移的坐标系里，位姿数值高精度但坐标系不共享，广播出来的数字没有公共基准。\textbf{正确理解}：异构融合的核心不是“融合估计值”，而是先“对齐坐标系”（跨越鸿沟）。没有共同坐标系，再准的位姿也拼不到一起——必须先配准（路线一/二）或改求相对位姿（路线三）。
\end{pitfall}

\begin{pitfall}[思维陷阱：不问任务需求，一律追求全局一致地图]
\textbf{错误做法}：默认所有协同都需要一个全局一致的共享地图，不惜代价去做全局对齐。\textbf{现象}：在异构鸿沟下死磕全局对齐、耗费大量工程，而任务其实只需“无人机告诉地面车目标在哪”这种相对关系，全局地图是过度设计。\textbf{根本原因}：把“全局一致地图”当成协同的必需品，没区分任务真正需要的一致性层级。\textbf{正确理解}：先问“我的协同任务到底需要全局一致还是相对一致？”引导/跟踪/相对避障类任务往往只需相对位姿（路线三，简单得多）；真正需要全局地图的（如多机协同建图）才上路线一/二。
\end{pitfall}

\begin{practice}
\begin{enumerate}
  \item （实现 $+$ 实验）用本节代码（或 Open3D 的 3D ICP）做 $2.5\mathrm{D}$ 配准实验：生成两个有重叠的 3D 点云（模拟无人机视觉点云和地面激光点云，故意让密度不同、覆盖部分重叠），先各自投影成高程图再用 ICP 配准。验证 (a) 配准后两坐标系对齐；(b) 对照——直接 ICP 原始点云，观察密度/覆盖差异如何让它失败。然后把点云压平成一个平面，观察 ICP 如何退化，理解路线一“需几何特征”的局限。
  \item （分析 $+$ 设计）对四个场景分别判断存在哪些鸿沟（视角/模态/坐标系）并选融合路线：(a) 两架无人机都用下视相机协同建图；(b) 无人机相机 $+$ 地面车激光、GPS 拒止室内搜救；(c) 无人机引导地面车前往一个无人机看到的目标（户外）；(d) 无人机热成像 $+$ 地面车 RGB 相机夜间协同。
  \item （跨章综合，串联 SLAM 与本章）路线一的 ICP、路线三的局部 BA 都是 SLAM 章的几何工具。说明：(a) 同构多机 SLAM 里这些工具怎么用（前提是什么）；(b) 异构融合里复用它们时输入发生了什么关键改变（从“原始观测”变成“精选的中间表示”）；(c) 论证“异构感知融合没有发明新的几何优化工具，而是发明了让异构观测变得可比的‘中间表示’和‘问题降级’（全局$\to$相对）”，并结合 \cref{sec:hmr-formation} 的“输出空间一致性”说明“找到合适的公共表示/公共空间”是不是贯穿整个异构协同的元方法。
\end{enumerate}
\end{practice}

% ════════════════════════════════════════════════════════════════
\section{异构协同操作：不对称 wrench cone 与内力分配}
\label{sec:hmr-manip}

\paragraph{动机：均分内力，吸盘却被要求“推”。}
协同操作的力学核心：$n$ 个机器人在物体上的接触点施加力/力矩，grasp matrix $\mathbf{G}$ 把这些接触 wrench 映射到物体合力旋量 $w_{\mathrm{obj}}=\mathbf{G}f$。要让物体按期望运动需 $\mathbf{G}f=w_{\mathrm{des}}$；这个方程通常欠定（接触自由度 $>$ 物体 $6$ 自由度），解空间是特解 $+$ $\mathbf{G}$ 零空间，零空间里的力是\textbf{内力（internal wrench）}——不产生物体净运动，但维持抓取稳定（如对抗性的“夹紧力”）。同构操作默认各 agent 力可行集对称，内力在各接触点\emph{均分}。现在把接触异构化，考虑一个异构联合抓取：三个 agent 共同搬一个箱子——左边是带夹爪的地面机械臂（夹爪能推能拉，双边摩擦锥），右边是带吸盘的机械臂（吸盘吸在箱子侧面，\emph{只能拉、不能推}——你没法用吸盘“推”一个面），上方是用绳索吊住箱子的无人机（绳索\emph{只能沿绳方向拉、且拉力必须为正}——绳子不能压）。均分内力策略立刻出问题：内力的本质是“对抗性的夹紧”，但吸盘\emph{不能推}（若均分要求它施加“推”的内力分量，它做不到），绳索吊的无人机\emph{只能沿绳拉、且拉力为正}（若均分要求横向力或压力，它也做不到）。均分隐含假设“每个 agent 能施加任意方向的力”，这只对对称的双边摩擦锥成立。\textbf{异构接触下，每个 agent 的力被限制在一个形状各异的可行锥里，均分会要求某些 agent 施加它们物理上无法施加的力。}

\paragraph{反面：如果沿用对称均分会怎样。}
沿用均分内力喂给异构接触，会发生：（1）\textbf{吸盘端要求负压（推力）$\to$ 物理不可实现}（均分算出的吸盘内力含指向物体的分量（推），但吸盘只能拉，控制器要么饱和、要么物体被另外两个 agent 推得偏离期望——抓取失稳）；（2）\textbf{绳索端要求侧向力或负张力 $\to$ 绳子做不到}（绳子只能沿自身方向拉，侧向分量无法提供，负张力意味着“绳子推物体”（松弛）、物体下坠）；（3）\textbf{物体净运动偏离期望}（因为某些 agent 无法施加分配给它们的力，实际合力 $\mathbf{G}f\ne w_{\mathrm{des}}$）。根本原因：\textbf{均分把内力当成“可以任意分配的对抗力”，但异构接触下每个 agent 的可施加力被锁在各自的 wrench cone 里，分配必须尊重这些不对称约束。}

\paragraph{历史：从对称抓取到异构联合操作。}
协同操作的力学（grasp matrix、内力、force closure）源自多指手抓取理论（Murray, Li, Sastry 的经典著作），那里假设各手指接触模型相同（点接触 $+$ 摩擦、软接触等），可行力锥对称；机器人协同搬运早期也沿用这个对称假设。异构操作的需求随空中操作和混合编队兴起：Michael 等\cite{michael2014aerial} 的多四旋翼绳索悬挂协同搬运是一个典型——绳索接触是单边的（只能拉），是处理\textbf{不对称接触约束}的代表性工作。% [待修] 原句称其“首次/标志性大规模处理不对称接触约束”，已删此优先权断言：见下方 \rebuilt——年份与“首次”均有误。
\rebuilt{两处更正：(1) 源标注“Auton. Robots 2014”有误——该文会议版为 RSS 2009、期刊版为 Auton. Robots 30(1):73--86, 2011（doi:10.1007/s10514-010-9205-0），无 2014 版，bib 键年份亦应改；(2) “协同操作里第一次大规模处理不对称(单边/只能拉)接触约束”属夸大：同组先有 Cheng/Fink/Kim/Kumar 平面协同拖曳(WAFR 2008 / ASME JMR 2009，含单边张力约束)及更早的抓取力封闭文献，该文真正新意是 3D 空中、绳索悬挂载荷的静平衡与稳定性分析，而非不对称约束的优先权。}Tognon 等（RA-L 2019）\cite{tognon2019aerial} 的空中冗余操作进一步处理飞行基座的力约束。近年的工作（如 2025 的力自适应近-远机器人协同操作、多人形搬运）则把异构接触推向更复杂的组合。这条线的核心进展是：\textbf{从“对称可行锥下的内力均分”走向“不对称可行锥下的内力优化”}——内力分配从一个解析的均分公式，变成一个带锥约束的优化问题。

\subsection{理论：在不对称 wrench cone 上做内力分配}
\label{sec:hmr-manip-qp}
为每个 agent $i$ 定义它在接触点能施加的力/力矩集合——wrench cone $\mathcal{W}_i$：
\begin{itemize}
  \item \textbf{夹爪（双边摩擦锥）}：标准摩擦锥 $\{f:\|f_t\|\le\mu f_n,\ f_n\ge 0\}$（法向可正可负取决于是否双边夹持，能推能拉）；
  \item \textbf{吸盘（单边 $+$ 拉限）}：只能拉，法向力 $f_n\le 0$（指向远离物体），切向受真空吸力限制 $\|f_t\|\le\mu|f_n|$；
  \item \textbf{绳索悬挂（射线）}：力只能沿绳方向 $\hat d$ 且张力为正，$f=\lambda\hat d,\ \lambda\ge 0$，这是最窄的可行集——一条射线。
\end{itemize}
这三个锥形状天差地别（实心圆锥 / 反向圆锥 / 一条射线），这正是异构的体现。物体期望旋量 $w_{\mathrm{des}}$，要求 $\mathbf{G}f=w_{\mathrm{des}}$（$f$ 是所有接触 wrench 的堆叠）。\textbf{内力分配问题}：
\begin{equation}
  \begin{aligned}
  \min_{f}\quad & f^\top W f &&(\text{最小化总力/能耗，或加权偏好})\\
  \mathrm{s.t.}\quad & \mathbf{G}f = w_{\mathrm{des}} &&(\text{合力满足期望})\\
  & f_i \in \mathcal{W}_i,\quad \forall i &&(\text{每个 agent 的力在各自可行锥内}).
  \end{aligned}
  \label{eq:hmr-internal-qp}
\end{equation}
对比同构操作：均分是 \cref{eq:hmr-internal-qp} 在“所有 $\mathcal{W}_i$ 相同且对称”时的解析特例（对称性让最优解恰好均分）。异构时 $\mathcal{W}_i$ 各异、不对称，QP 没有解析均分解，必须数值求解（摩擦锥用二阶锥约束 SOCP，或多面体近似后用 QP）。

\begin{insight}[“硬约束进可行集”原则在操作层的又一实例]
异构内力分配和 \cref{sec:hmr-cbba} 的能力硬约束、\cref{sec:hmr-formation} 的避障，本质上是同一个原则的不同实例——\textbf{“物理上绝对不可以”的事情（吸盘不能推、绳子不能压、机器人不会飞）必须编码进可行集（约束的形状），而不是代价。}内力分配里，这个原则体现为“力必须落在 wrench cone 内”是硬约束，写进 QP 的约束而非目标。这就是为什么异构操作的内力分配必然是带锥约束的优化，而不能是一个均分公式——均分公式无法表达“某些方向的力绝对不可施加”这件事。
\end{insight}

\paragraph{异构带来的可行性问题。}
异构接触还引入一个同构操作里少见的问题：\textbf{抓取可能不可行（infeasible）}。如果各 agent 的 wrench cone 太窄、组合起来无法张成 $w_{\mathrm{des}}$ 所需方向，QP 无解——物体根本无法被这组异构接触按期望操作。例如三个 agent 全是绳索悬挂（都只能拉），它们能合成的力都是“向上拉”的组合，\emph{无法提供任何向下的净力}——如果任务需要把物体往下压，这组抓取不可行。这就是为什么异构操作要检查 \textbf{force closure（力封闭）}：各 agent 的 wrench cone 的组合（Minkowski 和）必须张成整个 wrench 空间（能抵抗任意方向的外部扰动）。前置自测第 $4$ 题的答案在这里落地——异构接触的可行 wrench 空间是各 agent 可行集的 Minkowski 和。

\begin{insight}[对比性思维：内力不是对称分配的对抗力，而是锥约束下的优化解]
异构协同操作的内力\textbf{不是}“对称分配给每个 agent 的对抗力”，\textbf{而是}“在各 agent 形状各异的可行锥约束下、为满足合力且保持抓取稳定而求出的优化解”。前者假设力可任意分配（同构对称），后者承认每个 agent 的力被锁在不对称的可行集里。把异构操作当对称操作处理，会要求 agent 施加它们物理上不可能的力、导致抓取失稳——这是异构操作最典型的失败模式。
\end{insight}

\noindent 下面给出异构内力分配的 QP 骨架。为什么不能像同构那样用解析均分？均分能成立依赖一个隐含条件：所有接触的力可行集相同且对称（关于“推/拉”对称），对称性让“最优内力恰好均分”成为解析结论。异构接触下这个对称性消失——吸盘只能拉、绳子只能沿绳拉，可行集既不相同也不对称，此时“均分”会落到某些 agent 的可行集之外（要求吸盘推、绳子压），物理不可实现。所以必须显式把“每个 agent 的力 $\in$ 它的 wrench cone”写成约束、在这些约束下优化——这只能是带约束 QP/SOCP。

\begin{codebox}[Python]{异构内力分配的带锥约束 QP（多面体近似）}
import numpy as np
from scipy.optimize import minimize

def allocate_internal_force(G, w_des, agent_types, normals, W=None):
    """异构内力分配:min f'Wf s.t. Gf=w_des, f_i in wrench_cone_i.
       G: grasp matrix (6, sum_contact_dim);返回各接触力 f."""
    n_dim = G.shape[1]
    W = np.eye(n_dim) if W is None else W            # 默认最小化总力平方
    A_blocks, b_blocks, off = [], [], 0
    for atype, nrm in zip(agent_types, normals):     # 组装每个 agent 的锥约束（不等式）
        Ai, bi, dim = build_wrench_cone_constraints(atype, nrm)  # 夹爪/吸盘/绳索各异
        Af = np.zeros((Ai.shape[0], n_dim)); Af[:, off:off+dim] = Ai
        A_blocks.append(Af); b_blocks.append(bi); off += dim
    A_ineq, b_ineq = np.vstack(A_blocks), np.concatenate(b_blocks)
    cons = [
        {"type": "eq",   "fun": lambda f: G @ f - w_des},        # 合力满足期望（硬）
        {"type": "ineq", "fun": lambda f: b_ineq - A_ineq @ f},  # 力在锥内（硬）
    ]
    res = minimize(lambda f: 0.5*f@W@f, np.zeros(n_dim), jac=lambda f: W@f,
                   constraints=cons, method="SLSQP")
    if not res.success:                              # 关键:异构抓取可能不可行
        raise RuntimeError("内力分配不可行:异构 wrench cone 无法张成期望旋量，"
                           "检查 force closure（各锥的 Minkowski 和是否覆盖所需方向）")
    return res.x
\end{codebox}

\noindent 三类典型错误：\emph{沿用对称均分不管接触类型}（均分出的力可能落在吸盘/绳索可行集之外、要求吸盘推、绳子压——物理不可实现、控制器饱和、抓取失稳）；\emph{把锥约束当软惩罚加进目标}（软惩罚下优化器可能输出一个轻微违反锥的解，但物理上吸盘一点都推不了——力可行集是硬物理约束）；\emph{不检查可行性、假设总有解}（异构抓取可能根本不可行，不检查会返回无意义的失败解、物体操作失败且不知原因）。\cref{tab:hmr-manip-compare} 对比同构与异构内力分配。

\begin{table}[H]\centering\small
\caption{同构 vs 异构内力分配}
\label{tab:hmr-manip-compare}
\begin{tabular}{lll}
\toprule
维度 & 同构（对称摩擦锥） & 异构（不对称 wrench cone） \\
\midrule
内力分配 & 解析均分公式 & 带锥约束 QP/SOCP（数值） \\
可行集 & 各 agent 相同、对称 & 各 agent 形状各异、不对称 \\
约束类型 & 对称摩擦锥 & 夹爪锥 / 吸盘半锥 / 绳索射线 \\
可行性 & 一般可行（对称足够） & 可能不可行（需查 force closure） \\
“推/拉”对称 & 是 & 否（吸盘只拉、绳子只拉） \\
\bottomrule
\end{tabular}
\end{table}

\begin{pitfall}[编程陷阱：异构接触沿用对称内力均分]
\textbf{错误做法}：用均分公式分配内力、不区分夹爪/吸盘/绳索。\textbf{现象}：吸盘被要求施加推力、绳子被要求施加压力或侧向力——物理不可实现、控制器饱和、物体偏离期望甚至脱落。\textbf{根本原因}：均分公式假设各接触力可行集对称，异构接触（单边/射线）违反此假设。\textbf{正确做法}：为每个 agent 构造其 wrench cone、解带锥约束的 QP/SOCP。\textbf{自检}：分配后检查每个力是否落在对应锥内（$f_{\mathrm{suction}}\cdot n\le 0$、$f_{\mathrm{cable}}$ 平行绳向且张力 $>0$）。
\end{pitfall}

\begin{pitfall}[概念误区：“只要有足够多的机器人，异构抓取总能完成任意操作”]
\textbf{错误理解}：agent 越多能合成的力越全，所以异构抓取总可行。\textbf{现象}：三根绳索吊一个箱子，任务要求把箱子向下压一下，怎么都做不到。\textbf{根本原因}：可行性取决于各 wrench cone 的 Minkowski 和能否张成所需方向，不取决于 agent 数量——三根只能拉的绳子，加再多根也提供不了向下的净力。\textbf{正确理解}：异构抓取的可行性由 force closure 决定（各锥 Minkowski 和的覆盖性），必须主动检查；某些方向若所有 agent 都无法提供，再多 agent 也不可行。
\end{pitfall}

\begin{pitfall}[思维陷阱：把内力分配的不可行当成求解器 bug]
\textbf{错误做法}：QP 求解失败时怀疑求解器配置/初值问题、反复调参重试。\textbf{现象}：换求解器、调初值、放松收敛阈值都没用，QP 始终无解。\textbf{根本原因}：可能不是数值问题，而是问题本身不可行——异构 wrench cone 组合无法满足 $\mathbf{G}f=w_{\mathrm{des}}$，这是物理事实、不是求解器缺陷。\textbf{正确理解}：QP 不可行先检查 force closure（问题层面），而非纠结求解器（数值层面）；异构操作的“不可行”是常态而非异常，要在设计阶段就检查可行性。\textbf{自检}：单独验证“各锥的 Minkowski 和是否包含 $w_{\mathrm{des}}$ 方向”——若不包含，无论怎么调求解器都无解。
\end{pitfall}

\begin{practice}
\begin{enumerate}
  \item （实现 $+$ 实验）实现异构联合抓取的内力分配：$1$ 夹爪 $+$ $1$ 吸盘 $+$ $1$ 绳索悬挂搬一个箱子。给定期望旋量（水平移动 $+$ 保持姿态），求解内力分配并验证 (a) 每个力落在各自锥内；(b) 合力满足期望。然后构造一个不可行实例（要求一个所有 agent 都无法提供方向的力），验证可行性检查正确报错。对比：把同一问题用均分公式求解，检查均分结果违反了哪些锥约束。
  \item （推导题）对“三根绳索悬挂一个点质量”，写出每根绳的 wrench cone（射线），求它们的 Minkowski 和。证明这个和无法包含任何“向下的净力”方向（即无法把物体往下压），因此“向下压”任务对这组抓取不可行。然后说明：加入一个吸盘（吸在物体下表面）后可行 wrench 空间如何改变？
  \item （跨章综合，串联抓取与本章）本节“力进 wrench cone（硬约束）”、\cref{sec:hmr-cbba} 的“能力硬剔除”、\cref{sec:hmr-formation} 的“避障硬约束”都体现“物理不可能进可行集、代价高进目标函数”。请：(a) 用一句话概括这个统一原则；(b) 对每个场景指出什么是“物理不可能”（进约束）、什么是“代价高/有偏好”（进目标）；(c) 论证为什么把“物理不可能”错误地编码成“代价很大”在三个场景里都会导致同一类失败（优化器在边角情况下选中物理不可行解）。
\end{enumerate}
\end{practice}

% ════════════════════════════════════════════════════════════════
\section{异构 MARL：HAPPO 与异构动作空间}
\label{sec:hmr-marl}

\paragraph{动机：维度对不上，网络都搭不起来。}
前面几节用的是“模型驱动”的协同（分配、编队、操作都有显式模型）。但有些异构协同任务太复杂、交互太难建模，人们转向“学习驱动”——用多智能体强化学习（MARL）直接学策略。这立刻撞上一堵墙。最常用的协作 MARL 算法是 MAPPO（Multi-Agent PPO），它的一个关键工程实践是\textbf{参数共享（parameter sharing）}：所有 agent 共用一个策略网络 $\pi_\theta$，输入各自观测、输出各自动作。参数共享极大提升样本效率（所有 agent 的经验都用来训同一个网络），在\emph{同构}多智能体里几乎是标配。现在把 agent 异构化——一个“$2$ 无人机 $+$ $2$ 四足”的联合任务：无人机的动作是 $4$ 维（总推力 $+$ 三轴力矩，或归一化的姿态 $+$ 油门），四足的动作是 $12$ 维（每条腿 $3$ 个关节力矩）。你想用 MAPPO 的参数共享让一个网络同时输出两类动作——\textbf{网络的输出层维度该是 $4$ 还是 $12$？}$4$ 维喂不动四足（缺 $8$ 个关节指令），$12$ 维对无人机多了 $8$ 维（无意义）。动作空间维度不一致，参数共享的网络\emph{结构上无法定义}。即便用 padding（把无人机动作补零到 $12$ 维）勉强搭起网络，更深的问题暴露：无人机的第 $5$--$12$ 维动作是无意义的填充，但网络会试图学习它们、梯度被污染；而且两类 agent 的动作\emph{语义完全不同}（推力 vs 关节力矩），一个共享网络要同时学两套语义，训练几乎不收敛。这就是异构 MARL 的根本困难：\textbf{异构动作空间（维度 $+$ 语义不一致）让参数共享失效。}

\paragraph{反面：如果硬用参数共享会怎样。}
强行参数共享 $+$ padding 异构动作，会遇到：（1）\textbf{填充维度的梯度污染}（无人机被 padding 到 $12$ 维，第 $5$--$12$ 维是无效填充，但网络不知道它们无效、反向传播时这些维度也产生梯度、扰乱有效维度的学习）；（2）\textbf{语义冲突}（同一个网络的同一个输出维度，对无人机是“推力”、对四足是“某关节力矩”，网络被迫在一个输出上叠加两种物理语义、学到的是两者的混乱平均）；（3）\textbf{不收敛或退化}（训练曲线剧烈震荡或卡在低水平，因为网络在两套不兼容的动作语义间反复横跳）。根本原因：参数共享假设所有 agent 的策略可以用同一组参数表达，这要求它们的动作空间同构。出路是\textbf{放弃参数共享}——让每个（或每类）agent 有自己的策略网络，各自的输出维度匹配各自的动作空间。但这立刻带来新问题：多个独立更新的策略如何保证联合训练\emph{稳定、单调改进}？这正是 HAPPO 要解决的。

\paragraph{历史：从 MAPPO 到 HAPPO。}
MARL 的协作算法发展有两条路线：\emph{值分解路线}（QMIX 等，把联合 Q 分解为各 agent 局部 Q，但有单调性约束的结构限制）；\emph{策略梯度路线}（MAPPO 等，直接优化策略、参数共享提样本效率，在同构任务上是强基线）。但 MAPPO 缺乏\textbf{理论上的单调改进保证}（它把单智能体 PPO 的 trick 套到多智能体、理论基础不牢），且参数共享假设同构。Kuba 等（HAPPO/HATRPO, ICLR 2022\cite{kuba2022happo}）与 Zhong 等（HARL, JMLR 2024\cite{zhong2024harl}）\pz{源稿原作“JMLR 2022 / ICLR 2022”，年份/会议有误：HAPPO/HATRPO 发表于 ICLR 2022；不存在“JMLR 2022”版，真正的 JMLR 文章是 HARL，为 2024 卷 25。已就地更正venue字符串；定理与引理内容本身经核无误}提出的 \textbf{Heterogeneous-Agent Reinforcement Learning（HARL）}框架——核心算法 HAPPO（和 HATRPO）——填补了这两个空白：\emph{理论}上证明了\textbf{多智能体优势分解引理（multi-agent advantage decomposition lemma）}，给出多智能体策略改进的严格单调保证（这是 MAPPO 没有的）；\emph{异构}上通过\textbf{顺序更新（sequential update scheme）$+$ 独立策略（无参数共享）}，天然支持异构 agent。后续工作（HAMDPO、HAA2C、HADDPG、HATD3 等）把这套框架扩展到连续/离散动作、确定性策略等。本节聚焦 HAPPO 的两个核心机制：优势分解和顺序更新。

\subsection{理论：优势分解 \texorpdfstring{$+$}{+} 顺序更新}
\label{sec:hmr-marl-theory}
HAPPO 的精妙之处在于：它把“同时更新所有 agent 策略”这个困难问题，转化为“按顺序逐个更新、每个 agent 更新时把前面 agent 的更新纳入考虑”。

\paragraph{多智能体优势分解引理。}
核心理论工具。它把联合优势函数（joint advantage）分解为各 agent 的\textbf{顺序边际贡献}。对一个 agent 的排列 $i_1,i_2,\dots,i_n$，联合优势可以写成：
\begin{equation}
  A^{i_{1:n}}\big(s,\ a^{i_{1:n}}\big) = \sum_{m=1}^{n} A^{i_m}\big(s,\ a^{i_{1:m-1}},\ a^{i_m}\big).
  \label{eq:hmr-advantage-decomp}
\end{equation}
直观地说：联合动作带来的总优势，等于“逐个 agent 加入时各自的边际优势之和”——agent $i_m$ 的边际优势是在\emph{已知前面 $i_1,\dots,i_{m-1}$ 的动作}的条件下，它选择 $a^{i_m}$ 带来的额外优势。这和 CBBA 的“边际收益”、\cref{sec:hmr-cbba} 的联盟边际贡献是同一个“顺序边际”的思想——只不过这里是在策略空间里。

\begin{insight}[多视角理解：优势分解与异构 CBBA 的 bundle 构建结构同源]
优势分解和 \cref{sec:hmr-cbba} 异构 CBBA 的 bundle 构建在结构上高度相似——CBBA 里每个机器人按顺序把任务加入包、每个任务的价值是“在已有任务基础上的边际收益”；优势分解里每个 agent 按顺序更新策略、每个 agent 的改进是“在前面 agent 已更新基础上的边际优势”。相似之处在于：都用“顺序 $+$ 边际”把一个耦合的联合问题拆成可逐个处理的序列。但不同之处在于：CBBA 的顺序是为了\emph{分布式可解}（贪婪近似），优势分解的顺序是为了\emph{理论保证}（严格单调改进）——前者是工程近似，后者是数学定理。理解这个“顺序边际”思想在分配和学习里的双重出现，是把握异构协同的一条暗线。
\end{insight}

\paragraph{顺序更新方案。}
基于优势分解，HAPPO 的训练不是同时更新所有策略，而是\textbf{按顺序逐个更新}（\cref{alg:hmr-happo}）。
\begin{algo}[HAPPO 顺序更新]\label{alg:hmr-happo}
对每一轮训练：随机打乱 agent 顺序 $i_1,i_2,\dots,i_n$；设累积重要性比率 $M\leftarrow 1$；对 $m=1,\dots,n$：（i）用 $M$（前面 agent 更新带来的联合策略变化）修正 agent $i_m$ 的优势估计；（ii）用 PPO 的裁剪目标更新 agent $i_m$ 的\emph{独立}策略 $\pi^{i_m}$；（iii）把 agent $i_m$ 的策略变化（重要性比率）累乘进 $M$，传给下一个 agent。
\end{algo}
\noindent 关键点：\textbf{每个 agent 有独立策略 $\pi^{i_m}$}（No Parameter Sharing, NoPS）——这天然支持异构动作空间，每个 $\pi^{i_m}$ 的输出维度匹配自己的动作空间（无人机的网络输出 $4$ 维、四足的输出 $12$ 维，各自独立）；\textbf{顺序更新 $+$ 重要性修正}——每个 agent 更新时把前面 agent 已经做出的策略改变（通过累积重要性比率 $M$）纳入自己的优势估计，这保证了“我改进我的策略时，考虑了队友刚才的改进”，从而联合策略\emph{单调改进}。这两个机制配合，使得异构动作空间被自然支持（独立策略）、联合训练有严格的单调改进保证（优势分解 $+$ 顺序修正）——这正是 MAPPO（参数共享、无单调保证）做不到的。

\subsection{NoPS vs ParPS 的根本张力}
\label{sec:hmr-marl-tension}
这是本节最深刻、也最实用的洞察。HAPPO 的单调改进保证\textbf{依赖 NoPS（无参数共享）}——优势分解和顺序更新的理论推导假设每个 agent 有独立策略。\rebuilt{术语溯源更正：“NoPS / ParPS”这套命名与“ParPS$+$顺序更新$\Rightarrow$基线漂移$\Rightarrow$破坏单调保证”这条因果机制，出自 Yu 等 2025（OMDPG, arXiv:2507.09989, \cite{hmarl2025ompg}），\emph{并非} HARL 原文（\cite{zhong2024harl}）的表述——HARL 用的是 FuPS/SePS 术语，其消融只说“参数共享带来不合理的策略约束、损害单调改进（定理 7 假设异构性）”，未提“baseline drift”。本小节沿用 Yu 2025 的 NoPS/ParPS 框架来组织讨论是可以的，但勿把该框架与基线漂移机制回挂到 HARL/HAPPO 论文本身；归于 HARL 的只应是“无参数共享$+$顺序更新天然支持异构动作空间”这半句。}但实践中纯 NoPS 有一个代价：\textbf{样本效率下降}——每个 agent 单独学自己的策略、不能像参数共享那样“共享经验”，$n$ 个 agent 就要训 $n$ 个网络、每个只用自己的数据，样本效率远低于参数共享。而\textbf{部分参数共享（Partial Parameter Sharing, ParPS）}——让同类 agent 共享部分参数——能恢复一些样本效率。但近期研究（如 2025 的相关工作）\cite{hmarl2025ompg} 发现：直接把 ParPS 和 HAPPO 结合会引入\textbf{策略更新基线漂移（baseline drift）}问题——参数共享让 agent 的更新相互干扰、破坏了顺序更新赖以保证单调改进的独立性假设。

\begin{insight}[样本效率与单调改进保证不可兼得]
异构 MARL 面临一个深刻的张力——\textbf{样本效率（要参数共享）与单调改进保证（要独立策略）不可兼得}。参数共享让所有 agent 的经验汇集训练、样本效率高，但它假设策略同质、且破坏了 HAPPO 单调改进的理论前提；独立策略保证了单调改进和异构支持，但每个 agent 各学各的、样本效率低。这不是工程实现的缺陷，而是异构学习的\emph{内在权衡}——你想要异构灵活性 $+$ 理论保证，就得在样本效率上让步；你想要样本效率，就得牺牲一部分异构灵活性或理论保证。HAPPO 选择了前者（保证 $+$ 异构，牺牲效率），这个选择在样本昂贵的真实机器人上是有代价的——这也是为什么“UAV $+$ 四足的联合 HAPPO 训练至今没有大规模成功案例”\pz{联网未核实：此为否定存在性命题，原则上不可被普遍证实；对抗式检索（arXiv/Scholar，2024--2026）未找到反例，腿足多机 MARL 成功案例均为同构（清一色四足或双足），含 UAV$+$四足的 EMOS(ICLR'25) 用 LLM 高层规划$+$经典控制器、明确不用 MARL。措辞已设限（“至今/大规模”），与现有文献一致，保留旗标存疑而非当既成事实}。
\end{insight}

\paragraph{另一条思路：Transformer 统一异构输入。}
除了 HAPPO 的“独立策略”路线，处理异构动作空间还有另一条思路：用 \textbf{Transformer / attention 机制统一异构输入输出}。核心想法是把每个 agent 的观测/动作表示成\emph{变长的 token 序列}，用 attention 处理任意维度的输入，输出层用 agent 类型相关的 head 解码成各自的动作。这样一个模型能处理不同维度的 agent，又保留了一定的参数共享（attention 主干共享、解码 head 分类型）。两条思路的对比见 \cref{tab:hmr-marl-three}。

\begin{table}[H]\centering\small
\caption{异构 MARL 三种方案对比}
\label{tab:hmr-marl-three}
\begin{tabular}{llllll}
\toprule
方案 & 异构动作 & 参数共享 & 单调保证 & 样本效率 & 成熟度 \\
\midrule
MAPPO $+$ 共享 & 不支持 & 是 & 无 & 高 & 仅同构 \\
HAPPO（独立策略 NoPS） & 每 agent 独立网络 & 无 & 有（理论证明） & 低 & 较成熟 \\
Transformer 统一 & 变长 token $+$ 分类型 head & 部分（主干共享） & 无显式保证 & 较高 & 探索中 \\
\bottomrule
\end{tabular}
\end{table}

\begin{insight}[对比性思维：拥抱异构 vs 弥合异构]
处理异构动作空间，HAPPO 的策略\textbf{不是}“想办法把异构 agent 塞进一个共享网络”，\textbf{而是}“承认它们本质不同、给每个独立的策略，再用顺序更新协调它们的学习”。Transformer 路线反过来，\textbf{不是}“放弃共享”，\textbf{而是}“用足够灵活的架构（attention $+$ 分类型 head）让共享重新可能”。两条路线代表两种哲学：HAPPO 拥抱异构（独立 $+$ 协调），Transformer 弥合异构（灵活共享）。哪条更好至今没有定论——这正是异构 MARL 最活跃的前沿。
\end{insight}

\noindent 下面给出 HAPPO 顺序更新的核心逻辑骨架（伪实现，聚焦机制而非完整 RL 管线）。为什么不并行更新所有 agent、而要顺序更新？并行更新（所有 agent 同时改策略）的问题：每个 agent 基于“队友旧策略”计算自己的改进，但所有 agent 同时改了 $\to$ 实际环境是“所有队友都用新策略”、各自的改进假设失效 $\to$ 联合策略可能变差（无单调保证），这是 MAPPO 的隐患。顺序更新让 agent 逐个改、每个改的时候把“前面 agent 已经改了”纳入考虑（重要性修正 $M$），所以它的改进是基于“队友的真实当前策略”——这才能保证联合单调改进。

\begin{codebox}[Python]{HAPPO 顺序更新（独立策略 NoPS $+$ 累积比率修正优势）}
import numpy as np

class HeteroAgent:
    """异构 agent:独立策略网络，动作维度匹配自己的动作空间."""
    def __init__(self, obs_dim, act_dim, agent_type):
        self.act_dim = act_dim                  # UAV=4, 四足=12 —— 各自不同！
        self.agent_type = agent_type
        self.policy = build_policy_net(obs_dim, act_dim)   # 独立网络（NoPS）

def happo_sequential_update(agents, rollouts, clip=0.2, lr=3e-4):
    """随机顺序逐个更新独立策略，用累积比率 M 修正优势."""
    order = np.random.permutation(len(agents))  # 每轮随机打乱顺序（保公平与理论性质）
    M = np.ones(rollouts.batch_size)            # 累积联合策略变化（初值 1）
    old = [a.policy.clone() for a in agents]
    for m in order:
        agent = agents[m]; obs, act = rollouts.obs[m], rollouts.act[m]
        advantage = rollouts.joint_advantage * M               # 用 M 修正（关键）
        ratio = np.exp(agent.policy.logprob(obs, act) - old[m].logprob(obs, act))
        clipped = np.clip(ratio, 1 - clip, 1 + clip)
        loss = -np.minimum(ratio * advantage, clipped * advantage).mean()
        agent.policy.step(loss, lr)                            # 更新本 agent 独立策略
        new_ratio = np.exp(agent.policy.logprob(obs, act) - old[m].logprob(obs, act))
        M = M * new_ratio                                      # 累积，供后续 agent 修正
    return agents
\end{codebox}

\noindent 三类典型错误：\emph{参数共享 $+$ padding 异构动作}（总是输出最大维度 $12$、对 UAV 取前 $4$ 维、后 $8$ 维白算且产生梯度污染、两套语义在同一权重上冲突）；\emph{独立策略但并行更新}（丢掉 $M$ 修正——每个 agent 基于“队友旧策略”算改进、但所有人同时改了、联合策略实际变化未被各自考虑、无单调保证、训练可能发散）；\emph{顺序更新但用固定顺序}（固定顺序让排在前面的 agent 系统性占优、引入偏置，HAPPO 要求每轮随机打乱顺序）。\cref{tab:hmr-marl-choose} 给出方案选型。

\begin{table}[H]\centering\small
\caption{异构 MARL 方案选型}
\label{tab:hmr-marl-choose}
\begin{tabular}{ll}
\toprule
情形 & 推荐方案 \\
\midrule
同构 agent & MAPPO $+$ 参数共享（最省样本） \\
异构 agent $+$ 要单调改进保证 & HAPPO（NoPS，独立策略 $+$ 顺序更新；牺牲样本效率换保证） \\
异构 agent $+$ 样本宝贵 $+$ 可接受无保证 & Transformer 统一架构（部分共享、前沿探索） \\
\bottomrule
\end{tabular}
\end{table}

\begin{pitfall}[编程陷阱：异构动作空间用参数共享 $+$ padding]
\textbf{错误做法}：把所有 agent 动作 padding 到最大维度、共用一个策略网络。\textbf{现象}：训练剧烈震荡或卡在低水平不收敛、填充维度学到无意义的值。\textbf{根本原因}：参数共享假设动作空间同构；padding 引入无效维度（梯度污染）$+$ 不同 agent 的动作语义冲突（推力 vs 关节力矩共用权重）。\textbf{正确做法}：每个（或每类）agent 独立策略网络（NoPS）、输出维度匹配各自动作空间，用 HAPPO 的顺序更新协调独立策略的联合训练。\textbf{自检}：检查网络输出维度是否对每类 agent 都恰好匹配其动作空间（无填充）。
\end{pitfall}

\begin{pitfall}[概念误区：“HAPPO 只是支持异构的 MAPPO”]
\textbf{错误理解}：HAPPO $=$ MAPPO 加个异构支持、本质一样。\textbf{现象}：以为可以随便在 HAPPO 里加参数共享提效率，结果单调改进保证失效、训练不稳。\textbf{根本原因}：HAPPO 的核心不是“支持异构”，而是“优势分解 $+$ 顺序更新带来的单调改进保证”，这个保证严格依赖 NoPS（独立策略）——加参数共享会引入基线漂移、破坏保证\pz{“基线漂移”这一具体机制及术语出自 Yu 等 2025（\cite{hmarl2025ompg}），非 HARL 原文；详见 \cref{sec:hmr-marl-tension} 的溯源说明}。\textbf{正确理解}：HAPPO 和 MAPPO 是不同理论基础的算法——MAPPO 无单调保证、靠参数共享提效率；HAPPO 有单调保证、靠独立策略 $+$ 顺序更新。二者在“参数共享”上是对立的设计选择，不能简单混搭。
\end{pitfall}

\begin{pitfall}[思维陷阱：忽视异构 MARL 的样本效率代价，盲目上 HAPPO]
\textbf{错误做法}：看到“异构”就上 HAPPO、不评估样本预算。\textbf{现象}：在真实机器人（样本极其昂贵）上训 HAPPO，$n$ 个独立网络各学各的、样本需求是参数共享的数倍、训练周期长到不现实。\textbf{根本原因}：忽略了 NoPS 的样本效率代价——HAPPO 的单调保证和异构支持是用样本效率换的。\textbf{正确理解}：上 HAPPO 前先问“我有多少样本预算？”——样本充足（仿真海量采样）则 HAPPO 的效率劣势可接受；样本宝贵（真机）则考虑 Transformer 统一（部分共享）或先在仿真训好再迁移。异构 MARL 的方案选择必须把样本预算作为一等约束。
\end{pitfall}

\begin{practice}
\begin{enumerate}
  \item （实现 $+$ 实验）实现 HAPPO 顺序更新的核心逻辑（用 PettingZoo 的简单异构环境，或自定义“$2$ 个动作维度不同的 agent”的玩具环境）。验证 (a) 顺序更新 $+$ $M$ 修正时联合回报单调上升；(b) 对照——去掉 $M$ 修正（并行更新）时训练是否震荡/发散。消融：固定顺序 vs 随机打乱顺序，观察对训练稳定性和公平性的影响。
  \item （分析题）\cref{sec:hmr-cbba} 的异构 CBBA bundle 构建和本节 HAPPO 的顺序更新都用“顺序 $+$ 边际”处理耦合问题。详细对比：(a) 二者的“顺序”分别为了什么（分布式可解 vs 理论保证）；(b) 二者的“边际量”分别是什么（任务边际收益 vs 策略边际优势）；(c) CBBA 的顺序贪婪是近似（$50\%$ 界）、HAPPO 的顺序更新是精确单调改进——为什么一个近似一个精确（提示：离散组合 vs 连续策略空间，以及是否做了重要性修正）。
  \item （开放设计题，串联能力矩阵与异构 MARL）为“$2$ UAV（$4$ 维动作）$+$ $2$ 四足（$12$ 维动作）”的联合搬运设计 MARL 方案：(a) 用 \cref{sec:hmr-axes} 的能力互补矩阵说明为什么这个任务需要异构团队；(b) 在 HAPPO 和 Transformer 统一架构之间做选择，给出判据（样本预算、是否需要保证、工程成熟度）；(c) 讨论一个与 \cref{sec:hmr-perception} 的连接点——异构 agent 的观测也是异构的（UAV 鸟瞰、四足平视），你的方案如何处理异构观测（独立编码器？共享的跨模态编码器？）。
\end{enumerate}
\end{practice}

% ════════════════════════════════════════════════════════════════
\section{异构 vs 同构的本质边界与分层协同架构}
\label{sec:hmr-bound}

\paragraph{动机：异构不是同构的并集。}
前面七节逐层讲了异构如何改写协同的每一面。这一节把它们收口。整章反复强调一句话：\textbf{异构系统不是同构系统的简单拼接}。一个诱人的错觉是“异构系统无非是几种同构子系统放在一起，我把每种机型当一个同构小组分别处理，再拼起来就行”。这个想法在最表层（如各机型各自的底层控制）部分成立，但在协同的每一个真正耦合的层面都会崩溃——因为\textbf{协同的本质是跨个体的耦合，而异构让这种耦合发生在“不可直接比较的个体”之间}。下面沿五个层面（分配、感知、编队、操作、学习）逐一对比同构与异构，看异构在每一层引入了什么同构系统里根本不存在的新问题。

\subsection{五层对比总表}
\label{sec:hmr-bound-table}
\cref{tab:hmr-five-layer} 是全章最重要的一张表——它把前七节的核心区别浓缩到一处。

\begin{table}[H]\centering\footnotesize
\caption{异构 vs 同构的五层对比：每一层的“新问题”都是“个体不可比”的产物}
\label{tab:hmr-five-layer}
\begin{tabular}{p{1.5cm}p{2.6cm}p{3.3cm}p{3.0cm}p{0.9cm}}
\toprule
协同层面 & 同构系统怎么做 & 异构系统的\textbf{新问题} & 异构系统怎么解 & 对应节 \\
\midrule
任务分配 & 所有 agent 可竞标所有任务，区别仅在代价；CBBA 有 $50\%$ 界 & 能力约束（谁有资格）、联盟任务（需多机拼能力，超可加破坏 DMG） & 候选集硬剔除 $+$ 联盟竞标/分解 $+$ 能量代价；联盟部分界失效 & \cref{sec:hmr-cbba} \\
通信 & 同质节点、相近能力，拓扑由距离定 & 节点异构（通信半径、续航差异）；哪类节点当中继/leader & 中继/mesh/LF 按节点能力分布选；公共语义层 & \cref{sec:hmr-comm} \\
编队/运动 & 同一动力学模板，一致性写在状态空间 & 动力学结构异构（全向/非完整/欠驱动），状态空间不可比 & 一致性上移到公共输出空间；协调/执行分离；阻抗柔性连接 & \cref{sec:hmr-formation} \\
感知 & 视角/模态相近，大量共视，位姿广播可拼 & 视角鸿沟 $+$ 模态鸿沟，无共视，坐标系无法对齐 & $2.5\mathrm{D}$ 几何中间表示 / 跨视角学习描述子 / 相对定位解耦 & \cref{sec:hmr-perception} \\
操作 & 对称摩擦锥，内力均分（解析公式） & 接触异构（夹爪/吸盘/绳索），wrench cone 不对称，可能不可行 & 带锥约束 QP；force closure 可行性检查 & \cref{sec:hmr-manip} \\
学习 & 同构动作空间，参数共享，MAPPO & 异构动作空间（维度/语义不同），参数共享失效 & 独立策略 $+$ 顺序更新（HAPPO）；样本效率代价 & \cref{sec:hmr-marl} \\
\bottomrule
\end{tabular}
\end{table}

\begin{insight}[异构协同的元方法：为不可比的个体寻找公共空间]
通读 \cref{tab:hmr-five-layer}，一个统一的模式浮现——\textbf{异构在每一层引入的“新问题”，本质上都是“个体不可直接比较”导致的}。分配里能力不可比（要硬约束筛选）、编队里状态不可比（要投影到公共输出）、感知里观测不可比（要找几何中间表示）、操作里力可行集不可比（要在各自锥上优化）、学习里动作空间不可比（要独立策略）。同构系统的所有协同算法都隐含一个“个体可比”的前提（可平均、可相减、可共享参数），异构打破了这个前提。\textbf{异构协同的所有技术，归根结底都是在回答同一个问题：当个体不可直接比较时，如何找到一个让它们变得可比的“公共空间/公共表示”，并在那里协调。}
\end{insight}

\noindent 这个洞察给出贯穿全章的元方法——\textbf{寻找公共空间（finding the common ground）}，见 \cref{tab:hmr-common-ground}。

\begin{table}[H]\centering\small
\caption{“寻找公共空间”元方法：各层不可比的东西与找到的公共表示}
\label{tab:hmr-common-ground}
\begin{tabular}{lll}
\toprule
层面 & 不可比的东西 & 找到的公共空间/表示 \\
\midrule
分配 & 异构能力 & 能力向量空间（\cref{sec:hmr-axes} 的 trait 向量） \\
编队 & 异构状态 & 公共输出空间（世界系位置） \\
感知 & 异构观测 & 几何中间表示（高程图）/ 学习嵌入空间 \\
操作 & 异构力可行集 & 物体 wrench 空间（grasp matrix 的值域） \\
学习 & 异构动作 & 独立策略 $+$ 顺序协调（不强求统一动作空间） \\
通信 & 异构私有状态 & 公共语义层（共享话题） \\
\bottomrule
\end{tabular}
\end{table}

\subsection{分层协同架构：把五层组织起来}
\label{sec:hmr-bound-arch}
前面每层是分开讲的，但一个真实异构系统里它们必须协同工作。把它们组织成一个\textbf{分层协同架构}（\cref{fig:hmr-arch}）。

\begin{figure}[ht]
\centering
\begin{tikzpicture}[
  >=Stealth, font=\footnotesize, node distance=3mm,
  L/.style={draw, rounded corners, align=left, inner sep=4pt, text width=10.2cm},
  l4/.style={L, fill=main!8, draw=main!60!black},
  l3/.style={L, fill=second!8, draw=second!60!black},
  l2/.style={L, fill=gray!12, draw=gray!55!black},
  l1/.style={L, fill=third!10, draw=third!60!black},
  ar/.style={->, thick, gray!55!black}]
\node[l4] (l4) {\textbf{L4 任务层}：能力感知分配（\cref{sec:hmr-cbba}）\\
  输入 能力矩阵 $Q$、任务需求 $Y$ $\to$ 输出 谁做什么（含联盟）\quad\emph{公共空间：能力向量空间}};
\node[l3, below=of l4] (l3) {\textbf{L3 协调层}：编队几何 / 联合操作规划（\cref{sec:hmr-formation,sec:hmr-manip}）\\
  输入 任务目标、各 agent 当前输出 $\to$ 输出 每个 agent 的输出空间目标\quad\emph{公共空间：输出空间/物体 wrench 空间}};
\node[l2, below=of l3] (l2) {\textbf{L2 执行层}：各 agent 私有控制器（异构！）\\
  UAV 微分平坦 $+\SE(3)$ \,$|$\, 四足 WBC \,$|$\, 臂 逆运动学\quad（每类各说各话，\emph{无公共空间——也无需}）};
\node[l1, below=of l2] (l1) {\textbf{L1 感知层}：跨模态融合（\cref{sec:hmr-perception}）$+$ 通信（\cref{sec:hmr-comm}）\\
  输入 各 agent 异构观测 $\to$ 输出 一致世界模型 / 相对位姿\quad\emph{公共空间：几何中间表示/学习嵌入}};
\draw[ar] (l4)--(l3) node[midway,right,font=\scriptsize]{任务目标下发};
\draw[ar] (l3)--(l2) node[midway,right,font=\scriptsize]{输出目标下发};
\draw[ar] (l1.east) to[out=0,in=0] node[midway,right,font=\scriptsize]{状态/感知回传} (l2.east);
\end{tikzpicture}
\caption{异构系统的分层协同架构。核心是“上层找公共空间协调、下层各自私有执行”——L4/L3/L1 都在各自的公共空间里协调（能力空间、输出空间、几何表示），只有 L2 执行层是每类 agent 私有的、异构的、无需公共空间。}
\label{fig:hmr-arch}
\end{figure}

\noindent 这个架构的核心是\textbf{“上层找公共空间协调、下层各自私有执行”}——这正是 \cref{sec:hmr-formation} 的“协调/执行分离”原则推广到整个系统。

\begin{insight}[理论-工程桥接：架构的分层是个体不可比性的直接产物]
这个分层架构不是凭空设计的，而是被异构性“逼”出来的。同构系统里协调和执行可以在同一层（同一个状态空间），因为大家可比；异构系统里可比性只在精心选择的公共空间里成立，而执行必然在各自的私有空间，所以协调层和执行层必须分开。\textbf{架构的分层，是个体不可比性的直接产物}——这也解释了为什么异构系统普遍比同构系统架构更复杂（多了“投影到公共空间”和“从公共空间目标到私有执行”两道转换）。
\end{insight}

\paragraph{选型决策框架。}
把整章方法论收成一个决策框架，面对一个真实异构问题时按它走：\textbf{第 $1$ 步——沿三轴诊断异构性}（\cref{sec:hmr-axes}）：能力轴异构则任务分配要能力感知（\cref{sec:hmr-cbba}）；动力学轴结构异构则编队要输出空间一致性（\cref{sec:hmr-formation}）；感知轴有视角/模态鸿沟则需跨模态融合（\cref{sec:hmr-perception}）。\textbf{第 $2$ 步——判断任务耦合强度}：任务可单机完成（只是筛选 $+$ 加权）则异构 CBBA 安全（保 $50\%$ 界）；任务需多机拼能力（超可加联盟）则 CBBA 界失效、上联盟形成/联合优化；任务需紧密物理协同（合抬重物）则上协同搬运 $+$ \cref{sec:hmr-manip} 异构内力。\textbf{第 $3$ 步——判断是否需要学习}：任务可建模则优先模型驱动（可解释、无需训练）；任务难建模则上 MARL、异构动作空间用 HAPPO（\cref{sec:hmr-marl}）。\textbf{第 $4$ 步——选通信架构}（\cref{sec:hmr-comm}）：高可靠性优先则 mesh、大范围 $+$ 远端弱通信则中继、一强多弱 $+$ 结构稳定则 leader--follower。\textbf{第 $5$ 步——确定一致性层级}（\cref{sec:hmr-perception}）：需全局一致地图则 $2.5\mathrm{D}$ 配准/跨视角描述子、只需相对关系则相对定位解耦（更简单，优先考虑）。

\begin{pitfall}[概念误区：“把每种机型当同构小组分别处理再拼起来 $=$ 异构协同”]
\textbf{错误理解}：异构系统 $=$ 几个同构子系统的并集、分别处理各机型再拼接。\textbf{现象}：各机型的底层控制都正常，但一旦需要跨机型协同（联合搬运、共享地图、编队机动），系统就出问题——拼接处全是 bug。\textbf{根本原因}：协同的本质是跨个体耦合、异构让耦合发生在不可比个体间。“分别处理再拼”只在无耦合的底层（L2 执行）成立，在所有真正耦合的层面（分配/感知/编队/操作/学习）都需要专门的“公共空间”机制。\textbf{正确理解}：异构协同的难点恰恰在“拼接处”——即跨个体的协调层；必须为每个耦合层面找到让异构个体可比的公共空间（\cref{tab:hmr-five-layer}），而非假设拼接是平凡的。
\end{pitfall}

\begin{pitfall}[思维陷阱：在错误的层面处理异构性]
\textbf{错误做法}：试图在执行层（L2）统一异构（逼所有机型用同一控制器），或在协调层（L3/L4）保留私有差异（让分配器理解每类机型的底层动力学）。\textbf{现象（前者）}：用统一控制器逼非完整车做全向机动——做不到；\textbf{（后者）}：分配器代码里塞满各机型动力学细节、耦合爆炸、无法维护。\textbf{根本原因}：把“该统一的层”和“该分异的层”搞反了——异构性应在执行层（L2）保留（各自私有控制器），在协调层（L3/L4/L1）通过公共空间消解（投影到可比表示）。\textbf{正确理解}：分层架构的纪律是——协调层找公共空间求可比（消解异构），执行层各自私有（保留异构）。
\end{pitfall}

\begin{pitfall}[思维陷阱：过度异构化（同构问题用异构重武器）]
\textbf{错误做法}：明明是参数异构（同模型不同 $\max v$）的简单情形，却上输出空间一致性、HAPPO 独立策略等异构重武器。\textbf{现象}：系统复杂度暴涨、但收益甚微——因为问题本来用同构方法加增益调度就够了。\textbf{根本原因}：没区分“参数异构”（同构方法即可）和“结构异构”（才需异构方法），\cref{sec:hmr-axes} 的三轴诊断第一步就是判断每个轴是参数还是结构异构。\textbf{正确理解}：异构方法是有成本的（架构更复杂、样本效率更低、可行性更难保证）；只对“结构异构”的轴用异构方法，“参数异构”的轴用同构方法 $+$ 调度。先诊断，再按需施治，避免杀鸡用牛刀。
\end{pitfall}

\begin{practice}
\begin{enumerate}
  \item （综合分析题）给定三个系统，分别判断“是真异构还是伪异构”，并说明各轴该用同构还是异构方法：(a) $5$ 台同型号但最大速度不同的 AGV；(b) $3$ 台地面车 $+$ $2$ 架无人机的搜救队；(c) $10$ 台同型号无人机但分两组、装了不同传感器（相机 vs 激光）。对每个系统填写一张“三轴诊断表”（能力/动力学/感知各是参数异构还是结构异构、需不需要异构方法）。
  \item （架构设计题）为“四足 $+$ 无人机 $+$ 固定臂”的建筑检测系统画出完整的四层协同架构（L1--L4），标注每层的输入/输出、公共空间、以及哪些前面的工具可直接复用、哪些必须改写。特别说明：L2 执行层为什么不需要公共空间、而 L1/L3/L4 为什么需要。
  \item （跨章综合，全章收口）本章反复出现“为不可比的异构个体找公共空间”这个元方法。请：(a) 把六个层面（分配/通信/编队/感知/操作/学习）各自的“不可比的东西”和“找到的公共空间”列成表，用自己的话重述每一项为什么是公共空间；(b) 论证这个元方法和 \cref{sec:hmr-formation} 的“协调/执行分离”、本节的“分层架构”是同一思想的不同表述；(c) 反思：是否存在某种异构性，\emph{找不到}合适的公共空间？如果有，那类问题该怎么办？（开放性研究级问题，尝试构造一个例子并讨论。）
\end{enumerate}
\end{practice}

% ════════════════════════════════════════════════════════════════
\section{端到端案例：四足 \texorpdfstring{$+$}{+} 无人机联合搜救}
\label{sec:hmr-case}

\paragraph{场景设定。}
灾后建筑搜救：一栋部分倒塌的多层建筑，需要快速定位可能的幸存者。可用资源：\textbf{$1$ 架无人机（UAV）}（下视相机、鸟瞰感知范围大、续航 $20$ 分钟、载重极低、可飞越废墟）、\textbf{$2$ 台四足机器人（Quad）}（激光雷达 $+$ RGB、地面视角、越障能力强、续航 $2$ 小时、中等载重、可钻进倒塌区）、\textbf{$1$ 个地面基站}（计算和通信枢纽）。任务清单：（T1）全区航拍生成俯视地图（需飞行）；（T2）标记可疑热点区域；（T3）地面进入可疑区搜寻（需越障）；（T4）搬开挡路的轻型障碍（需载重，可能需两台四足联盟）。三轴异构如何在这个场景同时出现见 \cref{tab:hmr-case-axes}。

\begin{table}[H]\centering\small
\caption{三轴异构在搜救场景的体现与对应方法}
\label{tab:hmr-case-axes}
\begin{tabularx}{\linewidth}{lLL}
\toprule
轴 & 本场景的体现 & 用哪节的方法打通 \\
\midrule
能力异构 & UAV 能飞不能搬、Quad 能搬能越障不能飞；T4 可能需两台 Quad 联盟 & \cref{sec:hmr-axes} 能力矩阵 $+$ \cref{sec:hmr-cbba} 异构 CBBA \\
动力学异构 & UAV 欠驱动飞行、Quad 足式步态——可达集结构不同 & \cref{sec:hmr-formation} 输出空间协调 $+$ 各自执行层 \\
感知异构 & UAV 鸟瞰相机 vs Quad 平视激光——视角 $+$ 模态双鸿沟 & \cref{sec:hmr-perception} 跨模态融合（$2.5\mathrm{D}$ 配准） \\
\bottomrule
\end{tabularx}
\end{table}

\noindent 通信上，无人机续航短（$20\,\min$）、基站固定、四足续航长——这是“远端续航长节点 $+$ 短续航空中节点”的结构，\cref{sec:hmr-comm} 建议用\emph{混合架构}：四足当局部 leader（续航长、计算够），无人机当空中侦察 follower，基站经四足中继接入。

\paragraph{端到端管线。}
把全章方法串成一条管线：\emph{阶段 $0$（感知-L1）}——UAV 航拍 $\to$ 鸟瞰 3D 点云 $\to$ $2.5\mathrm{D}$ 高程图，Quad 激光 $\to$ 地面 3D 点云 $\to$ $2.5\mathrm{D}$ 高程图，ICP 配准两图建立 UAV$\leftrightarrow$Quad 相对变换 $\mathbf{T}$（坐标系对齐）；\emph{阶段 $1$（任务-L4）}——在对齐的地图上标记任务，能力矩阵 $Q$ $+$ 任务需求 $Y$ $\to$ 异构 CBBA，T1/T2 派 UAV、T3 派 Quad、T4 若需 $60\,\mathrm{kg}$ 则两台 Quad 联盟；\emph{阶段 $2$（协调-L3）}——UAV 飞航拍航线、Quad 编队进可疑区，输出空间协调 $+$ 阻抗柔性连接（Quad 绕地面障碍、UAV 不受影响）；\emph{阶段 $3$（执行-L2）}——UAV 微分平坦生成轨迹、Quad WBC 步态控制，各自私有控制器实现 L3 的输出目标；\emph{阶段 $4$（回传-L1）}——Quad 发现热点 $\to$ 经相对变换 $\mathbf{T}$ 把位置传给 UAV/基站 $\to$ 共享语义话题。下面给出把这条管线串起来的顶层编排骨架（调用前面各节实现的模块）。

\begin{codebox}[Python]{四足 $+$ 无人机联合搜救端到端编排}
import numpy as np

def heterogeneous_sar_pipeline(uav, quads, base_station, tasks):
    # === 阶段 0:跨模态感知融合（6.5）—— 对齐坐标系 ===
    uav_elev, _ = pointcloud_to_elevation(uav.aerial_survey())          # 鸟瞰 2.5D
    quad_elev, _ = pointcloud_to_elevation(np.vstack([q.lidar_scan() for q in quads]))
    R_rel, t_rel = icp_2d(elev_to_points(quad_elev), elev_to_points(uav_elev))
    frame = (R_rel, t_rel)                                              # UAV<->Quad 相对变换

    # === 阶段 1:能力感知任务分配（6.1 + 6.2）===
    robots = [uav] + quads
    Q = np.array([r.capability_vector() for r in robots])              # 能力矩阵（6.1）
    assignment = hetero_cbba(robots, tasks, Q,
                             np.array([t.demand_vector() for t in tasks]),
                             [t.required_dims() for t in tasks],
                             np.array([r.battery() for r in robots]),
                             energy_model=sar_energy_model, enable_coalition=True)  # T4 触发联盟

    # === 阶段 2+3:协调（6.4）+ 执行（L2）===
    for robot, assigned in assignment.items():
        for task in assigned:
            if robot.type == "uav":                                   # 输出目标 -> 微分平坦（私有执行）
                robot.fly_differential_flat(task.target_in_frame(frame, to="uav"))
            else:
                wp = task.target_in_frame(frame, to="quad")
                acc = impedance_target(robot.pos, robot.vel, wp, np.zeros(2),
                                       robot.local_ground_obstacles())  # 6.4 阻抗
                robot.walk_wbc(acc)
                if task.needs_coalition:                              # T4 两台联盟搬运
                    f = allocate_internal_force(task.grasp_matrix, task.w_des,
                                                task.agent_types, task.normals)  # 6.6
                    robot.apply_contact_force(f[robot.contact_idx])

    # === 阶段 4:跨模态回传（6.3 公共语义层）===
    for q in quads:
        if q.found_hotspot():
            base_station.publish("/shared/target_pose",
                                 transform_to_world(q.hotspot_pos(), frame))
\end{codebox}

\begin{insight}[整条管线就是“在各种公共空间之间转换 $+$ 在每个公共空间里协调”]
这段编排代码最值得注意的，不是它调用了多少模块，而是\textbf{每个模块之间传递的都是“公共空间里的量”}——阶段 $0$ 输出的 \texttt{frame} 是坐标系对齐（公共几何空间），阶段 $1$ 输入的是能力向量（公共能力空间），阶段 $2$ 传给执行层的是 \texttt{wp}/\allowbreak\texttt{acc}（公共输出空间目标），阶段 $4$ 发布的是世界系位置（公共语义）。\textbf{整条管线就是“在各种公共空间之间转换 $+$ 在每个公共空间里协调”}——这把 \cref{sec:hmr-bound} 的“寻找公共空间”元方法落到了最具体的代码层面。异构系统的工程，本质就是管理这些公共空间和它们之间的转换。
\end{insight}

\subsection{前沿工作与开放问题}
\label{sec:hmr-frontier}
把本章方法放到最新研究的坐标系里，看还有哪些没解决的硬骨头。

\begin{paper}[RoCo：LLM 驱动的异构协调]\label{paper:hmr-roco}
\textbf{问题}：异构机器人如何用语义而非数值匹配来分工。\textbf{核心思想}（Mandi 等, ICRA 2024\cite{mandi2024roco}）：把每个机器人委托给一个 LLM agent，让它们用自然语言“对话式”地分工——LLM 负责高层任务分解，底层用运动规划执行；按三个性质组织协作——子任务能否并行/串行、所有 agent 是否获得相同任务状态信息、执行时机器人间的邻近度。\textbf{关键结果}：% [待修] 原句“在 RoCoBench 上 GPT-4 达到 86.7% 成功率”系数值张冠李戴，已注释删除，正确事实见下方 \rebuilt。
\rebuilt{数值更正——$86.7\%$ 这个数确实出自该文，但\emph{不是} RoCoBench 基准成绩：它是 §3.2 一个独立的玩具实验（$5^3$ 三维栅格多智能体路径可行性探针，30 次运行、允许至多 5 次重试）的 GPT-4 成功率。RoCoBench 的真实结果（表 2，Dialog/GPT-4）是逐任务的（Pack Grocery 0.44、Arrange Cabinet 0.75、Sweep Floor 0.95、Make Sandwich 0.80、Sort Cubes 0.93、Move Rope 0.65），原文\emph{未给}聚合成功率，无任何单项等于 $86.7\%$。故“RoCoBench 上 86.7\%”属把玩具实验数误植为基准成绩并杜撰了不存在的聚合值，不应引用。}\textbf{与本书关系}：这给异构协调提供了一个全新的“语义层”——不再用能力向量的数值匹配（\cref{sec:hmr-cbba}），而用 LLM 对任务和能力的语义理解来分工。\textbf{开放问题}：LLM 的语义分工如何与底层物理约束（\cref{sec:hmr-manip} 的 wrench cone、\cref{sec:hmr-formation} 的可达集）严格对接？语义层的“看似合理”如何保证物理可行？
\end{paper}

\noindent \textbf{力自适应近-远协同操作}（2025）：针对一个机器人在远处、一个在近处协同操作的异构场景做力自适应协调，把 \cref{sec:hmr-manip} 的异构内力分配推向动态、自适应的方向——可行锥不仅形状各异，还随机器人相对位置变化。\textbf{多双足协同搬运的去中心化 MARL}（Pandit 等, CoRL 2024\cite{pandit2025biped}）：用去中心化策略让多个\emph{双足}（Cassie）机器人协同搬运载荷、且策略可在不同机器人数量/构型间泛化（无需重训），部分回应了 \cref{sec:hmr-marl} 的开放问题。% [待修] 原句称此为“异构形态（多人形）协调涌现”，系误述：该系统用的是【同构的同型双足 Cassie】刚性连到共载体，并非异构形态、也非“humanoid/人形”；其新意是对【相同双足的不同数量/构型】泛化。会议为 CoRL 2024(PMLR v270)、非 CoRL 2025。详见下方 \rebuilt。
\rebuilt{Pandit 等三处更正：(1) 会议是 CoRL 2024（第 8 届，Munich；PMLR v270, pandit25a），“CoRL 2025”是另一届（Seoul），bib 键年份亦应核；(2) “异构形态/多人形”属杜撰——平台是清一色同型双足 Cassie（无臂无头，论文自称 multi-biped 而非 humanoid），无任何形态异构；(3) 故它并不能作为“异构形态协调涌现”的例证，只能作为“同构双足、去中心化控制对机器人数目/构型泛化”的例证。}此外，“真正异构形态（UAV $+$ 四足）的联合 MARL 训练至今没有大规模成功案例”\pz{联网未核实：否定存在性命题不可普遍证实；对抗式检索（2024--2026）未见反例，详见 \cref{sec:hmr-marl-tension} 同一论断的存疑说明}——这是 \cref{sec:hmr-marl} 留下的、本领域最硬的开放问题之一。

\begin{insight}[对比性思维：研究前沿不是把同构方法做得更好，而是在新问题谱系上开疆拓土]
异构协同的研究前沿，\textbf{不是}“把同构方法做得更好”，\textbf{而是}“在同构系统里根本不存在的新问题上开疆拓土”——LLM 语义分工（同构系统不需要，因为大家能力一样）、跨模态融合（同构系统无鸿沟）、异构动作空间学习（同构系统参数共享即可）。这再次印证全章主线：\textbf{异构不是同构的难一点的版本，而是一个有自己独立问题谱系的领域。}
\end{insight}

\begin{practice}
\begin{enumerate}
  \item （综合实现题，累积项目）把本章各节实现的模块（\cref{sec:hmr-axes} 能力矩阵、\cref{sec:hmr-cbba} 异构 CBBA、\cref{sec:hmr-formation} 阻抗编队、\cref{sec:hmr-perception} $2.5\mathrm{D}$ 配准、\cref{sec:hmr-manip} 内力分配）串成本节的端到端管线，在一个简化仿真里跑通“无人机航拍 $\to$ 地面车进区搜寻”的最小闭环。验证三轴异构都被正确处理：(a) 任务正确分配；(b) 坐标系对齐；(c) 地面车绕开矮障碍而无人机不受影响。
  \item （设计 $+$ 批判题）阅读 RoCo 的核心思想，为本节搜救场景设计一个“LLM 分工层”替代 \cref{sec:hmr-cbba} 的异构 CBBA。\emph{批判性地}分析：(a) LLM 分工相比能力向量数值分配，优势在哪（语义灵活性）、风险在哪（可能给出物理不可行的分工）；(b) 如何用 \cref{sec:hmr-manip} 的 force closure 检查、\cref{sec:hmr-formation} 的可达集检查给 LLM 的分工“兜底”（拒绝物理不可行的方案并要求重新分工）。
  \item （开放研究题，全章终极综合）\cref{sec:hmr-marl} 指出“真正异构形态（UAV $+$ 四足）的联合 MARL 至今无大规模成功案例”。综合全章知识分析这个开放问题难在哪并提出研究方案草图：(a) 异构动作空间用 HAPPO 还是 Transformer；(b) 异构观测（鸟瞰 vs 平视）的编码器怎么设计；(c) 奖励如何设计才能体现“能力互补”而非各自为战；(d) 样本效率问题怎么缓解（仿真预训练 $+$ 真机微调？课程学习？）。
\end{enumerate}
\end{practice}

% ════════════════════════════════════════════════════════════════
\section*{本章常见误解汇总}
\label{sec:hmr-misconceptions}
\begin{table}[H]\centering\footnotesize
\caption{异构多机协同常见误解}
\label{tab:hmr-misconceptions}
\begin{tabular}{p{4.0cm}p{6.6cm}p{1.3cm}}
\toprule
误解 & 正确理解 & 出处 \\
\midrule
“异构就是型号不同” & 异构是能力/动力学/感知三轴上的差异，且每轴还分参数异构（同模型不同参数）和结构异构（不同模型）；型号标签太粗，丢失能力连续性和可累加性 & \cref{sec:hmr-axes} \\
“所有能力都能靠多机累加” & 能力分可累加（载重）和不可累加（最高高度/分辨率）；混淆会导致“两台 $60\,\mathrm{m}$ 无人机被派去飞 $100\,\mathrm{m}$” & \cref{sec:hmr-axes} \\
“机型越多，异构系统越强” & 异构红利来自能力覆盖性和互补性、不是机型数量；能力雷同的新机型只增冗余 & \cref{sec:hmr-axes} \\
“用大代价软惩罚表达‘做不了’” & 能力约束必须硬剔除（进可行集）、不能软惩罚（进代价）；软惩罚下不可行分配在边角情况会被选中 & \cref{sec:hmr-cbba} \\
“异构 CBBA 也有 $50\%$ 最优界” & 纯能力 $+$ 能量约束保持 $50\%$ 界；联盟任务（超可加）破坏 DMG、界失效 & \cref{sec:hmr-cbba} \\
“mesh 网络最鲁棒、所以总是最优” & 鲁棒性有代价（延迟、复杂度、带宽）；结构稳定/有天然 leader 时 LF 或中继更划算 & \cref{sec:hmr-comm} \\
“异构编队 $=$ 同构编队加增益调度” & 增益调度只处理参数异构；结构异构需把一致性从状态空间上移到输出空间，是架构改动 & \cref{sec:hmr-formation} \\
“异构融合 $=$ 把各自位姿广播拼接” & 核心不是融合估计值，而是先对齐坐标系（跨越视角/模态鸿沟）；无共同坐标系再准也拼不到一起 & \cref{sec:hmr-perception} \\
“手工描述子能跨视角匹配” & ORB/SIFT 的视角不变性远不足以跨越鸟瞰 vs 平视的 $90^\circ$ 鸿沟，需几何中间表示或学习描述子 & \cref{sec:hmr-perception} \\
“异构操作沿用对称内力均分” & 均分假设可行集对称；异构接触（吸盘/绳索单边）的 wrench cone 不对称，需带锥约束 QP；还可能不可行 & \cref{sec:hmr-manip} \\
“agent 够多，异构抓取总能完成任意操作” & 可行性由各 wrench cone 的 Minkowski 和覆盖性（force closure）决定、不是 agent 数量 & \cref{sec:hmr-manip} \\
“异构动作空间用参数共享 $+$ padding” & 参数共享假设动作空间同构；padding 引入梯度污染 $+$ 语义冲突，需独立策略（NoPS）$+$ 顺序更新 & \cref{sec:hmr-marl} \\
“HAPPO 只是支持异构的 MAPPO” & HAPPO 的核心是优势分解 $+$ 顺序更新带来的单调改进保证，严格依赖 NoPS；加参数共享会破坏保证 & \cref{sec:hmr-marl} \\
“把各机型当同构小组分别处理再拼 $=$ 异构协同” & 协同的本质是跨个体耦合，异构让耦合发生在不可比个体间；难点在拼接处（协调层），需为每层找公共空间 & \cref{sec:hmr-bound} \\
\bottomrule
\end{tabular}
\end{table}

% ════════════════════════════════════════════════════════════════
\section*{本章小结}
\label{sec:hmr-summary}
本章把同构协同工具箱系统地异构化。一条认知主线贯穿始终：\textbf{异构系统不是同构系统的简单拼接，因为协同的本质是跨个体耦合，而异构让耦合发生在“不可直接比较的个体”之间}。整章所有技术，归根结底都在回答同一个元问题——\textbf{当个体不可比时，如何找到让它们变得可比的“公共空间/公共表示”，并在那里协调}。
\begin{itemize}
  \item \cref{sec:hmr-axes} 用能力/动力学/感知三轴定义异构性，区分参数异构（同构方法即可）与结构异构（才需异构方法），并把能力写成 trait 向量、队伍写成能力矩阵，区分可累加/不可累加特质——全章的共同语言。
  \item \cref{sec:hmr-cbba} 把同构 CBBA 异构化（能力硬约束 $+$ 联盟任务 $+$ 能量约束），核心洞察是“物理不可能进可行集、代价高进目标”，以及联盟任务的超可加效用如何破坏 $50\%$ 界。
  \item \cref{sec:hmr-comm} 给出空地通信的中继/mesh/leader--follower 三架构及其单点故障谱，强调跨类型通信必须走公共语义层。
  \item \cref{sec:hmr-formation} 揭示异构编队的核心原则——一致性必须建在公共输出空间而非状态空间，协调/执行分离，阻抗柔性连接吸收可达集差异。
  \item \cref{sec:hmr-perception} 攻克异构感知的视角鸿沟 $+$ 模态鸿沟，给出 $2.5\mathrm{D}$ 几何配准、跨视角学习描述子、相对定位解耦三条路线，核心是“选对让异构观测可比的中间表示”。
  \item \cref{sec:hmr-manip} 把内力分配从对称均分推广到不对称 wrench cone 上的带约束优化，强调 force closure 可行性检查。
  \item \cref{sec:hmr-marl} 用 HAPPO 的优势分解 $+$ 顺序更新处理异构动作空间，揭示“样本效率 vs 单调改进保证”的根本张力（NoPS vs ParPS）。
  \item \cref{sec:hmr-bound} 用五层对比总表系统论证“异构 $\ne$ 同构并集”，提炼“寻找公共空间”元方法和分层协同架构。
  \item \cref{sec:hmr-case} 把全章方法组装成四足 $+$ 无人机搜救的端到端管线，落地累积项目。
\end{itemize}

\subsection*{术语速查表}
\begin{table}[H]\centering\footnotesize
\caption{本章术语速查}
\label{tab:hmr-glossary}
\begin{tabular}{p{4.6cm}p{6.4cm}p{1.0cm}}
\toprule
术语（中英） & 一句话定义 & 首见 \\
\midrule
异构性三轴（capability/dynamics/perception heterogeneity） & 沿能力、动力学、感知三个正交轴刻画机器人队伍的差异 & \cref{sec:hmr-axes} \\
参数异构 vs 结构异构 & 同一模型不同参数（增益调度可解）vs 不同模型/约束（需专门方法） & \cref{sec:hmr-axes} \\
能力/特质向量（trait vector） & 用连续数值向量刻画机器人在各能力维度上的强度 & \cref{sec:hmr-axes} \\
能力矩阵 $Q$ & 行 $=$ 机器人、列 $=$ 能力维度的稠密矩阵，异构分配的输入 & \cref{sec:hmr-axes} \\
能力互补矩阵 $C^{\mathrm{comp}}$ & 度量任意两机器人能力互补程度，越大越互补 & \cref{sec:hmr-axes} \\
可累加/不可累加特质（cumulative/non-cumulative trait） & 能否通过多机相加满足任务需求（载重可、最高高度不可） & \cref{sec:hmr-axes} \\
物种（species） & 按特质聚类的机型组，物种内同构、物种间异构 & \cref{sec:hmr-axes} \\
合格候选集 $\mathcal{A}_j$ & 能力达标、有资格竞标任务 $j$ 的机器人集合 & \cref{sec:hmr-cbba} \\
联盟任务（coalition task） & 单机能力不足、需多机能力拼合的任务 & \cref{sec:hmr-cbba} \\
DMG（边际收益递减） & 子模性，CBBA $50\%$ 界的前提 & \cref{sec:hmr-cbba} \\
超可加效用（super-additive utility） & 多一个合作者收益升高（联盟任务），破坏 DMG & \cref{sec:hmr-cbba} \\
AGHS（空地异构系统） & UAV$+$UGV 协作的空地异构系统 & \cref{sec:hmr-comm} \\
公共语义层（shared semantic layer） & 异构节点交换信息的公共表示（世界系位姿等） & \cref{sec:hmr-comm} \\
公共输出空间（shared output space） & 异构 agent 状态投影到的同维同义空间（世界系位置） & \cref{sec:hmr-formation} \\
协调/执行分离 & 在公共输出空间协调、在私有状态空间执行 & \cref{sec:hmr-formation} \\
阻抗引导（impedance-based guidance） & 用虚拟弹簧-阻尼柔性连接 follower，允许临时偏离队形 & \cref{sec:hmr-formation} \\
视角鸿沟（viewpoint gap） & 鸟瞰 vs 平视外观差异巨大，跨视角特征匹配失效 & \cref{sec:hmr-perception} \\
模态鸿沟（modality gap） & 不同传感器数据类型不同（图像 vs 点云），无法直接比对 & \cref{sec:hmr-perception} \\
$2.5\mathrm{D}$ 高程图（elevation map） & 2D 网格 $+$ 每格高度，视角无关的几何中间表示 & \cref{sec:hmr-perception} \\
相对定位解耦（semi-distributed relative localization） & 只求两两相对位姿、与各自状态估计解耦，回避全局对齐 & \cref{sec:hmr-perception} \\
wrench cone & 一个 agent 在接触点能施加的力/力矩可行集 & \cref{sec:hmr-manip} \\
force closure（力封闭） & 各 wrench cone 的 Minkowski 和能否张成整个 wrench 空间 & \cref{sec:hmr-manip} \\
HAPPO（Heterogeneous-Agent PPO） & 用优势分解 $+$ 顺序更新处理异构 agent 的 MARL 算法 & \cref{sec:hmr-marl} \\
多智能体优势分解 & 联合优势 $=$ 各 agent 顺序边际优势之和 & \cref{sec:hmr-marl} \\
顺序更新（sequential update） & 按随机顺序逐个更新策略、用累积比率修正优势 & \cref{sec:hmr-marl} \\
NoPS / ParPS & 无参数共享（HAPPO 需要）/ 部分参数共享（提效率但破坏保证） & \cref{sec:hmr-marl} \\
寻找公共空间（finding the common ground） & 为不可比的异构个体找到可比的公共表示——全章元方法 & \cref{sec:hmr-bound} \\
\bottomrule
\end{tabular}
\end{table}

\subsection*{知识点总表}
\begin{table}[H]\centering\footnotesize
\caption{本章知识点总表（难度：$\star$ 入门、$\star\star$ 进阶、$\star\star\star$ 硬核）}
\label{tab:hmr-knowledge}
\begin{tabular}{clp{5.4cm}cc}
\toprule
编号 & 知识点 & 核心要点 & 对应节 & 难度 \\
\midrule
6-1 & 异构性三轴定义 & 能力/动力学/感知；参数 vs 结构异构 & \cref{sec:hmr-axes} & $\star\star$ \\
6-2 & 能力矩阵与互补矩阵 & trait 向量、$Q$、$C^{\mathrm{comp}}$、可累加性 & \cref{sec:hmr-axes} & $\star\star$ \\
6-3 & 联盟能力聚合 & 可累加维 sum、不可累加维 max & \cref{sec:hmr-axes} & $\star\star$ \\
6-4 & 异构 CBBA 三改动 & 能力硬约束 $+$ 联盟任务 $+$ 能量代价 & \cref{sec:hmr-cbba} & $\star\star\star$ \\
6-5 & 硬约束 vs 软约束 & 物理不可能进可行集、代价高进目标 & \cref{sec:hmr-cbba} & $\star\star\star$ \\
6-6 & $50\%$ 界何时失效 & 联盟超可加破坏 DMG & \cref{sec:hmr-cbba} & $\star\star\star\star$ \\
6-7 & 三类通信架构 & 中继/mesh/LF 的单点故障谱 & \cref{sec:hmr-comm} & $\star\star$ \\
6-8 & 公共语义层 & 跨类型通信不传私有状态 & \cref{sec:hmr-comm} & $\star\star$ \\
6-9 & 输出空间一致性 & 异构动力学的一致性上移到输出空间 & \cref{sec:hmr-formation} & $\star\star\star$ \\
6-10 & 协调/执行分离 & 公共空间协调 $+$ 私有空间执行 & \cref{sec:hmr-formation} & $\star\star\star$ \\
6-11 & 阻抗引导 & 柔性连接吸收异构可达集差异 & \cref{sec:hmr-formation} & $\star\star\star$ \\
6-12 & 视角/模态鸿沟 & 异构感知融合的两道根本障碍 & \cref{sec:hmr-perception} & $\star\star\star\star$ \\
6-13 & 跨模态融合三路线 & $2.5\mathrm{D}$ 配准 / 学习描述子 / 相对定位 & \cref{sec:hmr-perception} & $\star\star\star\star$ \\
6-14 & 不对称内力分配 & 带 wrench cone 约束的 QP & \cref{sec:hmr-manip} & $\star\star\star$ \\
6-15 & force closure 可行性 & Minkowski 和覆盖性决定可行 & \cref{sec:hmr-manip} & $\star\star\star$ \\
6-16 & 异构动作空间 & 维度/语义不同，参数共享失效 & \cref{sec:hmr-marl} & $\star\star\star\star$ \\
6-17 & HAPPO 优势分解 $+$ 顺序更新 & 单调改进保证、依赖 NoPS & \cref{sec:hmr-marl} & $\star\star\star\star$ \\
6-18 & NoPS vs ParPS 张力 & 样本效率 vs 单调保证不可兼得 & \cref{sec:hmr-marl} & $\star\star\star\star$ \\
6-19 & 五层对比总表 & 异构 $\ne$ 同构并集 & \cref{sec:hmr-bound} & $\star\star\star$ \\
6-20 & 寻找公共空间元方法 & 全章统一思想 & \cref{sec:hmr-bound} & $\star\star\star$ \\
6-21 & 分层协同架构 & L1--L4，上层协调下层执行 & \cref{sec:hmr-bound} & $\star\star\star$ \\
6-22 & 端到端搜救管线 & 三轴异构在一个系统打通 & \cref{sec:hmr-case} & $\star\star$ \\
\bottomrule
\end{tabular}
\end{table}

\section*{延伸阅读}
\label{sec:hmr-reading}
\begin{itemize}
  \item \textbf{异构任务分配（能力感知）}：Gerkey \& Matari\'c, ``A Formal Analysis and Taxonomy of Task Allocation in Multi-Robot Systems,'' IJRR 2004\cite{gerkey2004taxonomy}（MRTA 分类法奠基）；Choi, Brunet, How, ``Consensus-Based Decentralized Auctions for Robust Task Allocation,'' IEEE T-RO 2009\cite{choi2009cbba}（CBBA 原始论文）；Ravichandar et al., ``STRATA,'' JAAMAS 2020\cite{ravichandar2019strata}\pz{源稿标注 ICRA 2019 有误，应为 JAAMAS 2020（亦见 AAMAS 2021）；arXiv:1903.05149 为 2019 预印本}（trait 向量、cumulative/non-cumulative、species 的来源）；长航时异构 MRTA（Calvo \& Capit\'an, arXiv:2411.02062, 2024; T-RO 2025）\cite{hmrta2024longendurance}（能量/续航约束、任务中继）。
  \item \textbf{空地协同（UAV-UGV）}：AGHS 综述（Zhou et al., \emph{Robotics and Autonomous Systems}, vol.\ 203, 2026; doi:10.1016/j.robot.2026.105514）\cite{aghs2024survey}\pz{源稿标注“2024、Machines/RAS”：经核期刊为 RAS（正确）、年份应为 2026（非 2024）、“Machines”一项不成立；bib 键年份 2024 亦应改}（全栈梳理）；Michael et al., ``Cooperative Manipulation and Transportation with Aerial Robots,'' Auton.\ Robots vol.\ 30(1):73--86, 2011（会议版 RSS 2009）\cite{michael2014aerial}\pz{源稿标注 2014 有误，无 2014 版}（空中协同搬运先驱）。
  \item \textbf{异构编队}：HetSwarm（2025）\cite{hetswarm2025}（阻抗引导的直接来源）。
  \item \textbf{跨模态融合}：$2.5\mathrm{D}$ 地图协同定位（Sensors 2017）\cite{aghs25dmap2017}（路线一代表）；半分布式跨模态相对定位（2025）\cite{crossmodal2025semidist}（路线三最新工作）。
  \item \textbf{异构 MARL}：Kuba et al., ``Trust Region Policy Optimisation in MARL (HAPPO/HATRPO),'' ICLR 2022\cite{kuba2022happo} 与 Zhong et al., ``Heterogeneous-Agent Reinforcement Learning,'' JMLR vol.\ 25, 2024\cite{zhong2024harl}\pz{源稿作“JMLR 2022 / ICLR 2022”有误：HAPPO/HATRPO=ICLR 2022，无 JMLR 2022；JMLR 文章是 HARL，2024}（核心，优势分解 $+$ 顺序更新）；异构 MARL 单调优化（Yu et al., OMDPG, arXiv:2507.09989; ICANN 2025）\cite{hmarl2025ompg}（NoPS vs ParPS 基线漂移——此术语与机制的真正出处）；Pandit et al., ``Learning Decentralized Multi-Biped Control for Payload Transport,'' CoRL 2024\cite{pandit2025biped}\pz{源稿作 CoRL 2025、且述为“异构形态”有误：会议为 CoRL 2024，平台为同构同型双足 Cassie，非异构形态}（同构双足、去中心化控制对机器人数目/构型泛化）。
  \item \textbf{LLM 驱动异构协调}：Mandi et al., ``RoCo: Dialectic Multi-Robot Collaboration with Large Language Models,'' ICRA 2024\cite{mandi2024roco}（LLM 对话式异构分工）。
\end{itemize}

\section*{本章与后续章节的关系}
\label{sec:hmr-relations}
\begin{table}[H]\centering\small
\caption{本章与后续章节的关系}
\label{tab:hmr-relations}
\begin{tabular}{lll}
\toprule
后续章节 & 关系 & 本章铺垫的知识点 \\
\midrule
双臂/多臂协同操作 & 异构操作的深化（不同臂、不同末端） & \cref{sec:hmr-manip} 异构 wrench cone、内力分配 \\
多足协同 loco-manipulation & 移动 $+$ 操作的异构协同 & \cref{sec:hmr-formation} 输出空间协调、\cref{sec:hmr-manip} 力分配 \\
人形 $+$ 四足 $+$ 臂混合体 & 最复杂的异构形态协同 & \cref{sec:hmr-axes} 三轴异构、\cref{sec:hmr-bound} 分层架构 \\
MARL 系列 & \cref{sec:hmr-marl} 异构 MARL 的系统展开 & HAPPO、优势分解、异构动作空间 \\
综合实战 & 整合本章异构层 $+$ MARL & 全章模块（累积项目异构层） \\
\bottomrule
\end{tabular}
\end{table}
\noindent \textbf{承上}：本章把同构协同工具（通信图、共识、CBBA、MPC 编队、grasp matrix）系统异构化，是从“同构理想”到“异构现实”的转折点。\textbf{启下}：后续章节在操作任务上深化异构、在学习上展开异构 MARL，本章是理解后续所有“真实异构系统”的地基。

\section*{故障排查手册}
\label{sec:hmr-troubleshooting}
\begin{table}[H]\centering\footnotesize
\caption{异构多机协同故障排查}
\label{tab:hmr-troubleshooting}
\begin{tabular}{p{3.4cm}p{3.4cm}p{4.2cm}c}
\toprule
症状 & 可能原因 & 排查步骤 & 相关节 \\
\midrule
物理上做不了的任务被派出去（无人机被派搬重物） & 能力约束用了大代价软惩罚而非硬剔除 & 检查是否维护合格候选集 $\mathcal{A}_j$；断言每个 $(i,j)$ 满足 $\mathbf{q}_i\succeq\mathbf{y}_j$；把软惩罚改为剔除 & \cref{sec:hmr-cbba} \\
联盟任务无人认领、静默丢失 & 需求超单机能力，\texttt{is\_eligible} 对所有单机返回 False，但未触发联盟逻辑 & 分配后检查未认领任务；若其需求超单机上限则是漏处理的联盟任务；启用联盟分解/竞标 & \cref{sec:hmr-cbba} \\
异构 CBBA 解远差于最优、无法解释 & 含联盟任务（超可加）破坏 DMG、$50\%$ 界失效 & 判断任务是“抢”（竞争）还是“凑”（合作）；合作部分换联盟形成算法 & \cref{sec:hmr-cbba} \\
无人机回去换电池时整个系统断联 & 用了中继型架构、中继无人机是单点故障且续航短 & 改 mesh 分散依赖；或多机轮班中继接力；检查枢纽节点续航 vs 任务时长 & \cref{sec:hmr-comm} \\
加一种新机型要改所有已有节点代码 & 跨类型通信传私有状态、无公共语义层 & 约定公共语义话题（世界系位姿）；每节点只发布投影后的公共语义；验证“加新机型改 $0$ 个旧节点” & \cref{sec:hmr-comm} \\
差速车在异构编队里原地打转/跟不上 & 在状态空间写刚性一致性、喂给非完整 agent & 把一致性上移到输出空间；协调层只给期望输出、执行层用非完整跟踪器；改阻抗柔性连接 & \cref{sec:hmr-formation} \\
异构编队遇障碍整队散开 & 用了刚性约束、某成员绕障破坏队形并干扰他人 & 改阻抗（柔性）连接；每个 agent 只对自己可达集相关障碍反应；调小阻抗刚度 $K_d$ & \cref{sec:hmr-formation} \\
各机器人自定位准、拼接后错位米级 & 跳过坐标系对齐、直接按各自 SLAM 坐标拼 & 做配准（$2.5\mathrm{D}$ ICP 或跨视角描述子）；或改求相对位姿；确认有共同坐标系基准 & \cref{sec:hmr-perception} \\
跨视角特征匹配全是错的 & 用手工描述子（ORB/SIFT）跨鸟瞰-平视、视角鸿沟太大 & 改 $2.5\mathrm{D}$ 几何中间表示 $+$ ICP；或用学习的跨视角描述子；可视化匹配线确认杂乱 & \cref{sec:hmr-perception} \\
ICP 配准不收敛 & 直接配原始点云（密度/覆盖差异）或场景无几何特征 & 先投影 $2.5\mathrm{D}$ 高程图再配；检查场景是否有高度变化；平坦区考虑加 fiducial & \cref{sec:hmr-perception} \\
异构抓取吸盘被要求推、物体失稳 & 沿用对称内力均分、落到吸盘可行集之外 & 为每 agent 构造 wrench cone；解带锥约束 QP；验证每个力落在对应锥内 & \cref{sec:hmr-manip} \\
内力分配 QP 反复求解失败 & 可能不是数值问题、是异构抓取本身不可行 & 先查 force closure（各锥 Minkowski 和是否覆盖 $w_{\mathrm{des}}$ 方向）；不可行则改抓取配置 & \cref{sec:hmr-manip} \\
异构 MARL 训练剧烈震荡/不收敛 & 参数共享 $+$ padding 异构动作（梯度污染 $+$ 语义冲突） & 改独立策略（NoPS）、输出维度匹配各自动作空间；用 HAPPO 顺序更新；检查无填充维度 & \cref{sec:hmr-marl} \\
HAPPO 加了参数共享后保证失效、不稳 & ParPS 引入基线漂移、破坏顺序更新的独立性假设 & 退回纯 NoPS；或用专门处理基线漂移的变体；接受样本效率代价 & \cref{sec:hmr-marl} \\
异构系统“拼接处”全是 bug & 把每机型当同构小组分别处理再拼、忽略跨个体耦合 & 识别哪些层有跨个体耦合；为每个耦合层找公共空间（五层总表）；按分层架构组织 & \cref{sec:hmr-bound} \\
\bottomrule
\end{tabular}
\end{table}
